% Mechanical Waves (Continued) — Physics Upper Sixth — Cours signé OnBuch+
% Official MINESEC syllabus. Compiler avec : tectonic PHY04-mechanical-waves-continued.tex
\newcommand{\DOCMATIERE}{Physics}
\newcommand{\DOCNIVEAU}{Upper Sixth}
\newcommand{\DOCMODULE}{Upper Sixth — Physics}
\newcommand{\DOCLECON}{4}
\newcommand{\DOCTITRE}{Mechanical Waves (Continued)}
% OnBuch+ — préambule « cours signés » ENRICHI (compilé avec tectonic / XeLaTeX)
% Système visuel « Pop » d'OnBuch+ : fond crème (jamais blanc), bordures encre
% épaisses, ombres « dures » décalées, titres Archivo Black en capitales.
% Version MATHÉMATIQUES : graphes (pgfplots/tikz), boîtes pédagogiques
% (définition, propriété, méthode, attention, le savais-tu, exercice résolu,
% exercice, TP, bilan), page de garde et signature OnBuch+ ancrée en bas.
%
% Macros attendues dans main.tex AVANT \input{preamble} :
%   \DOCMATIERE  \DOCNIVEAU  \DOCMODULE  \DOCLECON (numéro)  \DOCTITRE
\documentclass[11pt]{article}

\usepackage{fontspec}
\usepackage[english]{babel}
\usepackage[a4paper,margin=2cm,top=3.4cm,bottom=2.8cm,headheight=1.4cm,footskip=1.3cm]{geometry}
\usepackage{amsmath,amssymb,amsfonts}
\usepackage{mathtools} % \xmapsto, \coloneqq... souvent utilisés par les modèles
\usepackage{cancel} % simplifications barrées (bilans de forces)
% \abs{z} (module, valeur absolue) et \norm{u} (norme), utilisés par les modèles
\providecommand{\abs}{}\let\abs\relax\DeclarePairedDelimiter{\abs}{\lvert}{\rvert}
\providecommand{\norm}{}\let\norm\relax\DeclarePairedDelimiter{\norm}{\lVert}{\rVert}
\usepackage[table]{xcolor} % [table] : fournit aussi \rowcolors (alternance de lignes)
\usepackage{pagecolor}
\usepackage{tikz}
\usetikzlibrary{arrows.meta,positioning,calc,shapes.geometric,shapes.misc,shapes.arrows,decorations.pathmorphing,decorations.pathreplacing,decorations.markings,patterns,shadows,mindmap,trees,fit,backgrounds,matrix,intersections,angles,quotes}
\usepackage{pgfplots}
\usepackage[version=4]{mhchem} % équations nucléaires \ce{^{235}_{92}U}
\usepackage[european,siunitx]{circuitikz} % schémas électriques
\pgfplotsset{compat=1.18}
\usepgfplotslibrary{fillbetween,groupplots} % aires entre courbes (calcul intégral)
\usepackage{enumitem}
\usepackage{fancyhdr}
% [most] charge toutes les bibliothèques tcolorbox (listings, minted, raster,
% documentation...) et gonfle inutilement la mémoire TeX sur un document long
% (~300 boîtes) : on ne charge que ce qui est réellement utilisé.
\usepackage[skins,breakable]{tcolorbox}
\usepackage{graphicx}
\usepackage{colortbl}
\usepackage{booktabs}
\usepackage{tabularx}
\usepackage{diagbox} % case d'en-tête diagonale des tableaux à double entrée
% Colonnes extensibles centrée / gauche / droite (C, L, R), souvent utilisées par les modèles
\newcolumntype{C}{>{\centering\arraybackslash}X}
\newcolumntype{L}{>{\raggedright\arraybackslash}X}
\newcolumntype{R}{>{\raggedleft\arraybackslash}X}
\usepackage{array}
\usepackage{multicol}
\usepackage{siunitx}
% parse-numbers=false : accepte aussi \SI{2,5 \times 10^{-3}}{...}
\sisetup{locale=UK,output-decimal-marker={.},group-minimum-digits=4,parse-numbers=false}
\DeclareSIUnit\ohm{\ensuremath{\symup{\Omega}}} % U+2126 absent de la police
\DeclareSIUnit\division{div} % réglages d'oscilloscope (ms/div, V/div)
\DeclareSIUnit\tr{tr} % tours (vitesse de rotation, ex. \SI{1500}{\tr\per\minute})
\usepackage{qrcode}
% \degree : commande fréquemment utilisée par les modèles pour un angle ou une
% température en dehors de \SI{}{} (non fournie par les paquets chargés).
\providecommand{\degree}{^{\circ}}
% \qed (fin de démonstration), utilisé par les modèles sans amsthm
\providecommand{\qed}{\hfill$\square$}
% \barre : double barre des valeurs interdites dans un tableau de variations (tkz-tab)
\providecommand{\barre}{\ensuremath{\|}}
% environnement proof (amsthm non chargé) utilisé par les modèles
\ifdefined\proof\else\newenvironment{proof}[1][Proof]{\par\noindent\textit{#1.}\ }{\qed\par}\fi
\usepackage{lastpage}
\usepackage{hyperref}
\hypersetup{hidelinks}

% ============================================================
% Palette « Pop » OnBuch+ — mêmes hex que la classe P (plus_ui.dart)
% ============================================================
\definecolor{poporange}{HTML}{F59321}
\definecolor{popdark}{HTML}{E07A0C}
\definecolor{popcream}{HTML}{FFF6E8}
\definecolor{popink}{HTML}{1C1714}
\definecolor{poppurple}{HTML}{7A5AE0}
\definecolor{popgreen}{HTML}{2E9E6A}
\definecolor{poppink}{HTML}{E0457B}
\definecolor{popblue}{HTML}{2F7DE1}
\definecolor{popgold}{HTML}{C98A12}
\definecolor{popmuted}{HTML}{8A7F76}
\definecolor{popline}{HTML}{EADFD2}
\definecolor{popgray}{HTML}{8A7F76}
\colorlet{popgrey}{popgray}
% Alias tolérants (le modèle nomme parfois les couleurs différemment)
\colorlet{popwhite}{white}
\colorlet{popblack}{popink}
\colorlet{popred}{poppink}
\colorlet{popyellow}{popgold}
% Teintes claires pour fonds de tableaux / zones
\colorlet{popblueL}{popblue!10!white}
\colorlet{poporangeL}{poporange!14!white}
\colorlet{popgreenL}{popgreen!12!white}
\colorlet{poppurpleL}{poppurple!12!white}
\colorlet{poppinkL}{poppink!12!white}
\colorlet{popgoldL}{popgold!14!white}

\pagecolor{popcream}

% ============================================================
% Typographie
% ============================================================
\defaultfontfeatures{Ligatures=TeX}
\setmainfont{PlusJakartaSans-Medium.ttf}[Path=../../../fonts/,BoldFont=PlusJakartaSans-Bold.ttf]
% Maths en sans-serif (Fira Math) avec chiffres et lettres droites de Plus
% Jakarta, pour que les formules se fondent dans le texte.
\usepackage[math-style=ISO,bold-style=ISO]{unicode-math}
\setmathfont{FiraMath-Regular.otf}[Path=../../../fonts/]
\setmathfont{PlusJakartaSans-Medium.ttf}[Path=../../../fonts/,range={up/num,up/{Latin,latin},\mathup}]
% (icomma retiré : décimales avec point en anglais)
% Caractères mathématiques Unicode absents des polices de texte (Plus Jakarta,
% Archivo Black) : rendus par leur équivalent mathématique, en texte comme en
% formule — sinon ils s'affichent en carré vide (ex. « Arithmétique dans ℤ »).
\usepackage{newunicodechar}
\newunicodechar{ℝ}{\ensuremath{\mathbb{R}}}
\newunicodechar{ℤ}{\ensuremath{\mathbb{Z}}}
\newunicodechar{ℂ}{\ensuremath{\mathbb{C}}}
\newunicodechar{ℕ}{\ensuremath{\mathbb{N}}}
\newunicodechar{ℚ}{\ensuremath{\mathbb{Q}}}
\newunicodechar{𝔻}{\ensuremath{\mathbb{D}}}
\newunicodechar{α}{\ensuremath{\alpha}}
\newunicodechar{β}{\ensuremath{\beta}}
\newunicodechar{γ}{\ensuremath{\gamma}}
\newunicodechar{δ}{\ensuremath{\delta}}
\newunicodechar{ε}{\ensuremath{\varepsilon}}
\newunicodechar{θ}{\ensuremath{\theta}}
\newunicodechar{λ}{\ensuremath{\lambda}}
\newunicodechar{μ}{\ensuremath{\mu}}
\newunicodechar{σ}{\ensuremath{\sigma}}
\newunicodechar{τ}{\ensuremath{\tau}}
\newunicodechar{φ}{\ensuremath{\varphi}}
\newunicodechar{ψ}{\ensuremath{\psi}}
\newunicodechar{ω}{\ensuremath{\omega}}
\newunicodechar{Σ}{\ensuremath{\Sigma}}
% majuscules produites par \MakeUppercase dans les titres (α-aminés -> Α)
\newunicodechar{Α}{\ensuremath{\alpha}}
\newunicodechar{Β}{\ensuremath{\beta}}
\newunicodechar{Γ}{\ensuremath{\Gamma}}
\newunicodechar{Δ}{\ensuremath{\Delta}}
\newunicodechar{Ε}{\ensuremath{\varepsilon}}
\newunicodechar{Θ}{\ensuremath{\Theta}}
\newunicodechar{Λ}{\ensuremath{\Lambda}}
\newunicodechar{Μ}{\ensuremath{\mu}}
\newunicodechar{Τ}{\ensuremath{\tau}}
\newunicodechar{Φ}{\ensuremath{\Phi}}
\newunicodechar{Ψ}{\ensuremath{\Psi}}
\newunicodechar{Ω}{\ensuremath{\Omega}}
\newunicodechar{↦}{\ensuremath{\mapsto}}
\newunicodechar{∈}{\ensuremath{\in}}
\newunicodechar{∉}{\ensuremath{\notin}}
\newunicodechar{∧}{\ensuremath{\wedge}}
\newunicodechar{⇔}{\ensuremath{\Leftrightarrow}}
\newunicodechar{⟺}{\ensuremath{\Longleftrightarrow}}
\newunicodechar{⇒}{\ensuremath{\Rightarrow}}
\newunicodechar{⟹}{\ensuremath{\Longrightarrow}}
\newunicodechar{∘}{\ensuremath{\circ}}
\newunicodechar{∪}{\ensuremath{\cup}}
\newunicodechar{∩}{\ensuremath{\cap}}
\newunicodechar{≡}{\ensuremath{\equiv}}
\newunicodechar{⊂}{\ensuremath{\subset}}
\newunicodechar{⊥}{\ensuremath{\perp}}
\newunicodechar{∃}{\ensuremath{\exists}}
\newunicodechar{∀}{\ensuremath{\forall}}
\newunicodechar{∛}{\ensuremath{\sqrt[3]{\,}}}
\newunicodechar{∜}{\ensuremath{\sqrt[4]{\,}}}
\newunicodechar{⃗}{\ensuremath{{}^{\to}}}
\newunicodechar{ⁿ}{\ifmmode^{n}\else\textsuperscript{\ensuremath{n}}\fi}
\newunicodechar{ˣ}{\ifmmode^{x}\else\textsuperscript{\ensuremath{x}}\fi}
\newunicodechar{ᵇ}{\ifmmode^{b}\else\textsuperscript{\ensuremath{b}}\fi}
\newunicodechar{ʳ}{\ifmmode^{r}\else\textsuperscript{\ensuremath{r}}\fi}
\newunicodechar{ᵘ}{\ifmmode^{u}\else\textsuperscript{\ensuremath{u}}\fi}
\newunicodechar{ʸ}{\ifmmode^{y}\else\textsuperscript{\ensuremath{y}}\fi}
\newunicodechar{ᵃ}{\ifmmode^{a}\else\textsuperscript{\ensuremath{a}}\fi}
\newunicodechar{ᵉ}{\ifmmode^{e}\else\textsuperscript{\ensuremath{e}}\fi}
\newunicodechar{ᵏ}{\ifmmode^{k}\else\textsuperscript{\ensuremath{k}}\fi}
\newunicodechar{ᵖ}{\ifmmode^{p}\else\textsuperscript{\ensuremath{p}}\fi}
\newunicodechar{ᴺ}{\ifmmode^{N}\else\textsuperscript{\ensuremath{N}}\fi}
\newunicodechar{ᶜ}{\ifmmode^{c}\else\textsuperscript{\ensuremath{c}}\fi}
\newunicodechar{ʰ}{\ifmmode^{h}\else\textsuperscript{\ensuremath{h}}\fi}
\newunicodechar{ᵠ}{\ifmmode^{q}\else\textsuperscript{\ensuremath{q}}\fi}
\newunicodechar{ₙ}{\ifmmode_{n}\else\textsubscript{\ensuremath{n}}\fi}
\newunicodechar{ₐ}{\ifmmode_{a}\else\textsubscript{\ensuremath{a}}\fi}
\newunicodechar{ᵢ}{\ifmmode_{i}\else\textsubscript{\ensuremath{i}}\fi}
\newunicodechar{ᵣ}{\ifmmode_{r}\else\textsubscript{\ensuremath{r}}\fi}
\newunicodechar{ₖ}{\ifmmode_{k}\else\textsubscript{\ensuremath{k}}\fi}
\newunicodechar{ᵦ}{\ifmmode_{\beta}\else\textsubscript{\ensuremath{\beta}}\fi}
\newunicodechar{ₓ}{\ifmmode_{x}\else\textsubscript{\ensuremath{x}}\fi}
\newfontfamily\popdisplay{ArchivoBlack-Regular.ttf}[Path=../../../fonts/]
\newfontfamily\poplogo{SpaceGrotesk-Bold.ttf}[Path=../../../fonts/]
\newfontfamily\popsemi{PlusJakartaSans-SemiBold.ttf}[Path=../../../fonts/]
\newfontfamily\popxbold{PlusJakartaSans-ExtraBold.ttf}[Path=../../../fonts/]
\newfontfamily\popxboldls{PlusJakartaSans-ExtraBold.ttf}[Path=../../../fonts/,LetterSpace=3.0]

\setlist[enumerate]{leftmargin=1.7em,itemsep=5pt,topsep=4pt}
\setlist[itemize]{leftmargin=1.7em,itemsep=4pt,topsep=4pt,label={\color{poporange}\textbullet}}
\setlength{\parindent}{0pt}
\setlength{\parskip}{5pt}
\renewcommand{\arraystretch}{1.25}

% pgfplots : style Pop par défaut
\pgfplotsset{
  every axis/.append style={axis line style={popink,line width=1pt},tick style={popink},
    grid style={popline,line width=0.5pt},label style={font=\small},tick label style={font=\footnotesize}},
  popcurve/.style={poporange,line width=1.8pt,no markers,smooth},
}
% Même style disponible dans un \draw TikZ simple (hors pgfplots)
\tikzset{popcurve/.style={poporange,line width=1.8pt}}
\tikzset{popgrid/.style={popline,very thin}}
\colorlet{popcurve}{poporange} % le nom du style est parfois utilisé comme couleur

% ============================================================
% Pill / squiggle
% ============================================================
\newcommand{\poppill}[3]{%
  \tikz[baseline=-0.55ex]\node[draw=popink,line width=1.2pt,rounded corners=2.6pt,%
    fill=#2,inner xsep=6.5pt,inner ysep=3pt]{%
    \popxboldls\fontsize{7}{7}\selectfont\color{#3}\MakeUppercase{#1}};%
}
\newcommand{\popsquiggle}[1]{%
  \tikz[baseline=-0.4ex]{
    \draw[#1,line width=2.6pt,line cap=round]
      plot [smooth] coordinates {(0,0.09) (0.34,0.01) (0.68,0.21) (1.02,0.01) (1.36,0.10)};
  }%
}
% Mot-clé surligné (marqueur orange sous le texte)
\newcommand{\cle}[1]{{\popxbold\color{popdark}#1}}

% ============================================================
% En-tête / pied
% ============================================================
\pagestyle{fancy}
\fancyhf{}
\renewcommand{\headrulewidth}{0pt}
\renewcommand{\footrulewidth}{0pt}
\lhead{\raisebox{-0.15ex}{\poplogo\fontsize{13}{13}\selectfont\color{popink}OnBuch\color{poporange}+}}
\rhead{\poppill{\DOCMATIERE\ \textbullet\ \DOCNIVEAU\ \textbullet\ Chapter \DOCLECON}{white}{popink}}
\lfoot{\popxbold\fontsize{7.6}{7.6}\selectfont\color{popmuted}\MakeUppercase{Signed course OnBuch+}\ \textbullet\ {\popsemi\DOCTITRE}}
\rfoot{\tikz[baseline=-0.6ex]\node[rounded corners=8.5pt,fill=popink,minimum height=17pt,minimum width=17pt,inner xsep=5pt,inner sep=0]{\popxbold\fontsize{8}{8}\selectfont\color{popcream}\thepage\,/\,\pageref*{LastPage}};}

% ============================================================
% Titres
% ============================================================
\newcommand{\chaptitre}[2]{%
  \begin{tcolorbox}[enhanced,colback=poporange,colframe=popink,boxrule=2.6pt,arc=11pt,
      drop shadow=popink,left=15pt,right=15pt,top=12pt,bottom=11pt,before skip=3pt,after skip=18pt]
    \poppill{#2}{popink}{popcream}\par\vspace{9pt}
    {\raggedright\hyphenpenalty=10000\exhyphenpenalty=10000\popdisplay\fontsize{22}{24}\selectfont\color{popcream}\MakeUppercase{#1}\par}
  \end{tcolorbox}
}
% Section numérotée : pastille orange avec le numéro + titre Archivo
\newcounter{popsec}
\newcommand{\coursec}[1]{%
  \stepcounter{popsec}\setcounter{popsub}{0}%
  \par\vspace{16pt}\noindent
  \begin{minipage}{\linewidth}
  \tikz[baseline=-0.7ex]\node[draw=popink,line width=1.6pt,fill=poporange,rounded corners=4pt,minimum size=22pt,inner sep=0]{\popdisplay\fontsize{12}{12}\selectfont\color{popcream}\thepopsec};%
  \hspace{8pt}{\popdisplay\fontsize{15}{17}\selectfont\color{popink}\MakeUppercase{#1}}\par
  \vspace{4pt}\noindent{\color{poporange}\rule{36pt}{3.6pt}}
  \end{minipage}\par\nopagebreak\vspace{8pt}
}
\newcounter{popsub}
\newcommand{\courssub}[1]{%
  \stepcounter{popsub}%
  \par\vspace{9pt}\noindent{\popxbold\fontsize{11.5}{13}\selectfont\color{popink}{\color{poporange}\thepopsec.\thepopsub}\hspace{6pt}#1}\par\nopagebreak\vspace{4pt}
}

% ============================================================
% Boîtes pédagogiques — même recette : fond blanc, bordure épaisse colorée,
% ombre dure encre, badge plein. Titre optionnel [#1].
% ============================================================
\tcbset{popbase/.style={breakable,enhanced,colback=white,boxrule=2.3pt,arc=9pt,
  shadow={3pt}{-3pt}{0pt}{popink},left=14pt,right=14pt,top=11pt,bottom=11pt,before skip=11pt,after skip=15pt,
  fontupper=\popsemi\fontsize{10.6}{14.5}\selectfont\color{popink},
  fontlower=\popsemi\fontsize{10.6}{14.5}\selectfont\color{popink}}}
% Coche dessinée (le glyphe ✓ n'existe pas dans les polices du document)
\newcommand{\popcheck}[1]{\tikz[baseline=-0.5ex]\draw[#1,line width=1.6pt,line cap=round,line join=round](-0.13,0.0)--(-0.04,-0.09)--(0.13,0.1);}
\providecommand{\coche}{\popcheck{popgreen}}
\providecommand{\resultat}[1]{\par\smallskip\noindent\fcolorbox{popgreen}{popgreenL}{\parbox{\dimexpr\linewidth-2\fboxsep-2\fboxrule\relax}{\textbf{Result.} #1}}\par\smallskip}
\newcommand{\popboxhead}[3]{\poppill{#1}{#2}{white}\ifx\relax#3\relax\else\hspace{6pt}{\popxbold\fontsize{10.6}{12}\selectfont\color{popink}#3}\fi\par\vspace{7pt}}

\newtcolorbox{aretenir}[1][]{popbase,colframe=popgreen,before upper={\popboxhead{Key points}{popgreen}{#1}}}
\newtcolorbox{exemplebox}[1][]{popbase,colframe=poporange,before upper={\popboxhead{Exemple}{poporange}{#1}}}
\newtcolorbox{definition}[1][]{popbase,colframe=popblue,before upper={\popboxhead{Definition}{popblue}{#1}}}
\newtcolorbox{propriete}[1][]{popbase,colframe=poppurple,before upper={\popboxhead{Property}{poppurple}{#1}}}
\newtcolorbox{methode}[1][]{popbase,colframe=popgold,before upper={\popboxhead{Method}{popgold}{#1}}}
\newtcolorbox{attention}[1][]{popbase,colframe=poppink,before upper={\popboxhead{Warning}{poppink}{#1}}}
\newtcolorbox{experience}[1][]{popbase,colframe=popblue,colback=popblueL,before upper={\popboxhead{Experiment}{popblue}{#1}}}
% « Le savais-tu ? » : carte violette pleine (contexte, culture, Cameroun)
\newtcolorbox{savaistu}[1][]{popbase,colback=poppurple,colframe=popink,
  fontupper=\popsemi\fontsize{10.4}{14.2}\selectfont\color{white},
  before upper={\poppill{Did you know?}{white}{poppurple}\ifx\relax#1\relax\else\hspace{6pt}{\popxbold\color{white}#1}\fi\par\vspace{7pt}}}
% Exercice résolu : énoncé en haut, \tcblower puis la solution sur fond crème
\newcounter{popexo}\newcounter{popexr}
\newtcolorbox{exoresolu}[1][]{popbase,colframe=poporange,colbacklower=popcream,
  segmentation style={popink,line width=1.2pt,dashed},
  before upper={\stepcounter{popexr}\popboxhead{Worked example \thepopexr}{poporange}{#1}},
  before lower={\poppill{Solution}{popink}{popcream}\par\vspace{6pt}}}
% Exercice (non résolu en place) — la correction est dans la partie Corrigés
\newtcolorbox{exercice}[1][]{popbase,colframe=popink,
  before upper={\stepcounter{popexo}\popboxhead{Exercise \thepopexo}{popink}{#1}}}
% Corrigé
\newtcolorbox{corrige}[1][]{popbase,colframe=popgreen,colback=popgreenL,
  before upper={\popboxhead{Solution}{popgreen}{#1}}}
% Objectifs / prérequis / bilan
\newtcolorbox{objectifs}{popbase,colframe=popink,colback=poporangeL,
  before upper={\popboxhead{Chapter objectives}{popink}{}}}
\newtcolorbox{prerequis}{popbase,colframe=popink,colback=popblueL,
  before upper={\popboxhead{Prerequisites}{popblue}{}}}
\newtcolorbox{bilan}{popbase,colframe=popink,colback=white,boxrule=2.8pt,
  before upper={\popboxhead{Fiche bilan}{poporange}{L'essentiel à connaître pour l'examen}}}
% Figure encadrée avec légende
\newtcolorbox{popfigure}[1][]{enhanced,breakable=false,colback=white,colframe=popline,boxrule=1.4pt,arc=8pt,
  left=10pt,right=10pt,top=10pt,bottom=8pt,before skip=10pt,after skip=12pt,halign=center,
  lower separated=false,halign lower=center,
  fontlower=\popsemi\fontsize{9}{11.5}\selectfont\color{popmuted},
  #1}
\newcommand{\legende}[1]{\par\vspace{6pt}{\popsemi\fontsize{9}{11.5}\selectfont\color{popmuted}#1}}

% ============================================================
% Signature OnBuch+
% ============================================================
\newcommand{\poplogoinline}[1]{{\poplogo\fontsize{#1}{#1}\selectfont\color{popink}OnBuch\color{poporange}+}}
\newcommand{\signatureonbuch}{%
  \begin{tcolorbox}[enhanced,colback=popink,colframe=popink,boxrule=2.2pt,arc=10pt,
      drop shadow=poporange,left=14pt,right=14pt,top=10pt,bottom=10pt,before skip=0pt,after skip=0pt]
    \tikz[baseline=-0.7ex]\node[circle,fill=popgreen,draw=popcream,line width=1.3pt,minimum size=22pt,inner sep=0]{\popcheck{white}};%
    \hspace{9pt}\begin{minipage}[c]{0.62\linewidth}
      {\popxbold\fontsize{12}{13}\selectfont\color{popcream}Signed\ \textbullet\ {\poplogo OnBuch\color{poporange}+}}\par\vspace{2pt}
      {\raggedright\popsemi\fontsize{8}{10.5}\selectfont\color{popline}\MakeUppercase{OnBuch+ teaching team}\\\MakeUppercase{Follows the official MINESEC syllabus}\par}
    \end{minipage}\hfill
    {\poplogo\fontsize{22}{22}\selectfont\color{popcream}OnBuch\color{poporange}+}
  \end{tcolorbox}%
}

% ============================================================
% Page de garde
% #1 titre · #2 matière · #3 niveau · #4 module · #5 n° leçon · #6 durée ·
% #7 description / objectifs (texte ou itemize)
% ============================================================
\newcommand{\pagegarde}[7]{%
  \thispagestyle{empty}
  \newgeometry{a4paper,margin=1.7cm,top=1.6cm,bottom=1.4cm}%
  \begin{tikzpicture}[remember picture,overlay]
    \fill[poporange!18] ([xshift=-3.2cm,yshift=-3.6cm]current page.north east) circle (4.6cm);
    \fill[poppurple!14] ([xshift=2.4cm,yshift=7.6cm]current page.south west) circle (3.8cm);
  \end{tikzpicture}%
  {\poplogo\fontsize{30}{30}\selectfont\color{popink}OnBuch\color{poporange}+}\hfill
  \raisebox{0.6ex}{\poppill{Signed course \textbullet\ Premium}{popink}{popcream}}\par
  \vspace{1pt}\popsquiggle{poppurple}\par
  \vspace{8pt}
  \begin{tcolorbox}[enhanced,colback=poporange,colframe=popink,boxrule=2.8pt,arc=13pt,
      drop shadow=popink,left=18pt,right=18pt,top=12pt,bottom=13pt,after skip=10pt]
    \poppill{#2 \textbullet\ #3}{popink}{popcream}\hspace{5pt}\poppill{#4}{white}{popink}\par\vspace{8pt}
    {\popdisplay\fontsize{14}{14}\selectfont\color{popink}CHAPTER #5}\par\vspace{4pt}
    {\raggedright\hyphenpenalty=10000\exhyphenpenalty=10000\popdisplay\fontsize{25}{27}\selectfont\color{popcream}\MakeUppercase{#1}\par}
    \vspace{8pt}
    {\popxbold\fontsize{9.5}{9.5}\selectfont\color{popink}\MakeUppercase{Suggested time: #6}}
  \end{tcolorbox}
  \begin{tcolorbox}[enhanced,colback=white,colframe=popink,boxrule=2.2pt,arc=10pt,
      drop shadow=popink,left=16pt,right=16pt,top=10pt,bottom=10pt,after skip=8pt]
    \poppill{In this chapter}{poppurple}{white}\par\vspace{5pt}
    \popsemi\fontsize{9.3}{12.6}\selectfont\color{popink}#7
  \end{tcolorbox}
  \noindent\begin{minipage}[t]{0.487\linewidth}
    \begin{tcolorbox}[enhanced,colback=popgreen,colframe=popink,boxrule=2.2pt,arc=10pt,
        drop shadow=popink,left=11pt,right=11pt,top=10pt,bottom=10pt,height=3.1cm,valign=center]
      \tcbox[colback=white,colframe=popink,boxrule=1.3pt,arc=5pt,boxsep=0pt,nobeforeafter,
        left=3pt,right=3pt,top=3pt,bottom=3pt,valign=center]{\qrcode[height=1.9cm]{https://play.google.com/store/apps/details?id=com.luvvix.onbuch&hl=en}}%
      \hspace{8pt}\begin{minipage}[c]{0.5\linewidth}
        {\popdisplay\fontsize{11}{12.5}\selectfont\color{white}\MakeUppercase{Download OnBuch}}\par\vspace{4pt}
        {\raggedright\popsemi\fontsize{8.4}{11}\selectfont\color{white}Free \textbullet\ Google Play\\AI tutor, past papers, results\par}
      \end{minipage}
    \end{tcolorbox}
  \end{minipage}\hfill
  \begin{minipage}[t]{0.487\linewidth}
    \begin{tcolorbox}[enhanced,colback=poppurple,colframe=popink,boxrule=2.2pt,arc=10pt,
        drop shadow=popink,left=11pt,right=11pt,top=10pt,bottom=10pt,height=3.1cm,valign=center]
      \tikz[baseline=-0.7ex]\node[circle,fill=white,draw=popink,line width=1.3pt,minimum size=1.6cm,inner sep=0]{\popdisplay\fontsize{24}{24}\selectfont\color{poppurple}+};%
      \hspace{8pt}\begin{minipage}[c]{0.6\linewidth}
        {\popdisplay\fontsize{11}{12.5}\selectfont\color{white}\MakeUppercase{Go OnBuch+}}\par\vspace{4pt}
        {\raggedright\popsemi\fontsize{8.4}{11}\selectfont\color{white}All signed courses, unlimited AI tutor, PDF sheets — OnBuch+ tab in the app\par}
      \end{minipage}
    \end{tcolorbox}
  \end{minipage}
  \vspace{8pt}
  \popcontenu
  \vfill
  \signatureonbuch
  \restoregeometry
  \newpage
}

% Rangée « Ce cours contient » (page de garde)
\newcommand{\poptile}[3]{% #1 couleur #2 gros texte #3 légende
  \begin{tcolorbox}[enhanced,colback=#1,colframe=popink,boxrule=2pt,arc=9pt,shadow={2.5pt}{-2.5pt}{0pt}{popink},
      left=6pt,right=6pt,top=9pt,bottom=9pt,halign=center,valign=center,height=1.75cm,width=\linewidth,nobeforeafter]
    {\popdisplay\fontsize{12.5}{13}\selectfont\color{white}\MakeUppercase{#2}}\par\vspace{3pt}
    {\centering\popsemi\fontsize{7.8}{9.5}\selectfont\color{white}#3\par}
  \end{tcolorbox}}
\newcommand{\popcontenu}{%
  \par\noindent{\popxboldls\fontsize{7.5}{7.5}\selectfont\color{popmuted}THIS SIGNED COURSE CONTAINS}\par\vspace{7pt}
  \noindent\begin{minipage}[t]{0.235\linewidth}\poptile{popblue}{Course}{complete, illustrated, step by step}\end{minipage}\hfill
  \begin{minipage}[t]{0.235\linewidth}\poptile{poporange}{Methods}{and worked examples}\end{minipage}\hfill
  \begin{minipage}[t]{0.235\linewidth}\poptile{poppink}{Exercises}{exam style + solutions}\end{minipage}\hfill
  \begin{minipage}[t]{0.235\linewidth}\poptile{popgreen}{Summary sheet}{the essentials to remember}\end{minipage}\par}

% Sommaire
\newtcolorbox{sommaire}{enhanced,colback=white,colframe=popink,boxrule=2.2pt,arc=10pt,
  drop shadow=popink,left=18pt,right=18pt,top=13pt,bottom=13pt,before skip=6pt,after skip=20pt,
  before upper={\poppill{Contents}{poppurple}{white}\par\vspace{9pt}\popsemi\fontsize{10.6}{14.5}\selectfont\color{popink}}}

% Signature de fin de cours — poussée tout en bas de la dernière page
\newcommand{\findecours}{%
  \par\vfill\nopagebreak
  \begin{center}\popsquiggle{poporange}\end{center}\vspace{4pt}
  \signatureonbuch
}
\AtBeginDocument{\def\llbracket{\mathopen{[\mkern-3.5mu[}}\def\rrbracket{\mathclose{]\mkern-3.5mu]}}} % intervalles d'entiers (glyphe ⟦ absent de la police)
% Glyphes absents de Fira Math : redéfinis à partir de glyphes disponibles
\AtBeginDocument{%
  \def\setminus{\mathbin{\backslash}}\let\smallsetminus\setminus
  \def\vdots{\mathord{\vbox{\baselineskip3.2pt\lineskiplimit0pt\kern4pt\hbox{.}\hbox{.}\hbox{.}}}}%
  \def\ddots{\mathinner{\mkern1mu\raise6.5pt\hbox{.}\mkern2mu\raise3.6pt\hbox{.}\mkern2mu\raise0.7pt\hbox{.}\mkern1mu}}%
  \def\triangle{\mathord{\scalebox{0.9}{$\symup{\Delta}$}}}%
  \def\nabla{\mathord{\raisebox{\depth}{\scalebox{1}[-1]{$\symup{\Delta}$}}}}%
  \def\checkmark{\popcheck{popdark}}%
  \def\star{\mathbin{\ast}}\def\bigstar{\mathord{\ast}}%
  \def\diamond{\mathbin{\scalebox{0.6}{$\Diamond$}}}%
  \def\bigcup{\mathop{\mathchoice{\raisebox{-0.3em}{\scalebox{1.7}{$\cup$}}}{\scalebox{1.25}{$\cup$}}{\cup}{\cup}}\displaylimits}%
  \def\bigcap{\mathop{\mathchoice{\raisebox{-0.3em}{\scalebox{1.7}{$\cap$}}}{\scalebox{1.25}{$\cap$}}{\cap}{\cap}}\displaylimits}%
  \def\lesseqqgtr{\mathrel{\lessgtr}}\def\bigoplus{\mathop{\oplus}\displaylimits}%
}
\providecommand{\ssi}{if and only if} % abréviation fréquente
\AtBeginDocument{\providecommand{\wideparen}{\overparen}} % arc de cercle

%% FIGFIX-BEGIN v5 — correctifs globaux des figures (docs/onbuch-generation/tools/figfix)
\usepackage{adjustbox}\usepackage{etoolbox}\tcbuselibrary{raster}
% toute figure TikZ/pgfplots plus large que la ligne est réduite à la largeur disponible
\BeforeBeginEnvironment{tikzpicture}{\begin{adjustbox}{max width=\linewidth}}
\AfterEndEnvironment{tikzpicture}{\end{adjustbox}}
% légendes pgfplots lisibles par-dessus les courbes
\pgfplotsset{every axis/.append style={legend style={fill=white,fill opacity=0.92,text opacity=1,draw=black!15,inner sep=3pt}}}
% étiquettes d'arêtes (syntaxe edge["..."]) avec fond blanc
\tikzset{every edge quotes/.append style={fill=white,inner sep=1.2pt}}
%% FIGFIX-END
\begin{document}

\pagegarde{Mechanical Waves (Continued)}{Physics}{Upper Sixth}{Upper Sixth — Physics}{4}{2 periods}{This chapter completes the study of mechanical waves by examining two key topics: the polarization of transverse waves, which proves their transverse nature and has major technological applications, and the factors that determine the speed of transverse waves on taut strings and wires, essential for understanding musical instruments and wave transmission.
\vspace{4pt}
\begin{multicols}{2}\small
\begin{itemize}
\item By the end of this chapter I can define polarization and distinguish polarized from unpolarized transverse waves
\item By the end of this chapter I can demonstrate polarization using a mechanical model (rope and slits) and explain why longitudinal waves cannot be polarized
\item By the end of this chapter I can describe how plane polarized light is produced using a polarizer and detected with an analyser
\item By the end of this chapter I can state and apply Malus's law for the intensity of light transmitted through two polaroids (GCE extension)
\item By the end of this chapter I can list practical applications of polarization (sunglasses, LCD screens, stress analysis)
\item By the end of this chapter I can state the factors affecting the speed of a transverse wave on a taut string
\end{itemize}\end{multicols}}

\begin{sommaire}
\begin{multicols}{2}
\begin{enumerate}
\item Polarization of Transverse Waves: Concept and Mechanical Demonstration
\item Production and Detection of Plane Polarized Light; Applications
\item Speed of Transverse Waves on a Taut String: Factors and Formula
\item Experimental Study and Applications: The Sonometer and Vibrating Strings
\item Integration activity
\item Methods and worked examples
\item Exercises
\item Solutions to the exercises
\item Summary sheet
\end{enumerate}
\end{multicols}
\end{sommaire}

\begin{objectifs}
\begin{itemize}
\item By the end of this chapter I can define polarization and distinguish polarized from unpolarized transverse waves
\item By the end of this chapter I can demonstrate polarization using a mechanical model (rope and slits) and explain why longitudinal waves cannot be polarized
\item By the end of this chapter I can describe how plane polarized light is produced using a polarizer and detected with an analyser
\item By the end of this chapter I can state and apply Malus's law for the intensity of light transmitted through two polaroids (GCE extension)
\item By the end of this chapter I can list practical applications of polarization (sunglasses, LCD screens, stress analysis)
\item By the end of this chapter I can state the factors affecting the speed of a transverse wave on a taut string
\item By the end of this chapter I can derive and use the relation v = √(T/μ) where μ is the mass per unit length
\item By the end of this chapter I can combine v = fλ with v = √(T/μ) to solve problems on vibrating strings and musical notes
\item By the end of this chapter I can describe an experiment (sonometer) to investigate how wave speed depends on tension and linear density
\end{itemize}
\end{objectifs}

\begin{prerequis}
\begin{itemize}
\item Definition and properties of transverse and longitudinal waves (Lower Sixth / PHY03)
\item The wave equation v = fλ and the relation T = 1/f
\item Concepts of amplitude, wavelength, frequency and phase of a progressive wave
\item Basic notion that light is an electromagnetic (transverse) wave
\item Force, tension and mass per unit length (linear density) from mechanics
\end{itemize}
\end{prerequis}

\newpage

\coursec{Polarization of Transverse Waves: Concept and Mechanical Demonstration}

At a busy junction in Douala, a motorbike rider tilts his polarized sunglasses and suddenly the blinding glare from a car's windscreen disappears --- yet his friend with ordinary glasses still squints. Meanwhile, in a village near Bamenda, a musician tightens the string of his mvet and the note rises higher. What do sunglasses and a tightened string have in common? Both are governed by the behaviour of \cle{transverse waves}: one by the \emph{orientation} of their vibrations, the other by the \emph{speed} at which disturbances travel along a stretched medium. This chapter reveals the physics behind these everyday observations. We begin with the first idea: the orientation of vibrations, which leads us to the phenomenon of \cle{polarization}.

\courssub{Recall: Transverse and Longitudinal Waves}

Before we can speak of polarization, we must be absolutely clear about the two fundamental ways in which a mechanical disturbance can travel through a medium. The classification depends on one single question: \emph{in which direction do the particles of the medium vibrate, compared with the direction in which the wave travels?}

\begin{definition}[Transverse and longitudinal waves]
A wave is \cle{transverse} when the particles of the medium vibrate \textbf{perpendicular} to the direction of propagation of the wave. Examples: a wave on a stretched rope, ripples on the surface of water (approximately), electromagnetic waves.

A wave is \cle{longitudinal} when the particles of the medium vibrate \textbf{parallel} to the direction of propagation, producing alternating compressions and rarefactions. Example: sound in air.
\end{definition}

Picture a rope tied at one end to a wall in the school laboratory. If you jerk the free end up and down, a hump travels along the rope: each bit of rope moves \emph{vertically} while the hump moves \emph{horizontally}. That is a transverse wave. Now picture a long spring (a ``slinky'') lying on a bench: if you push and pull one end along the axis of the spring, coils bunch together and spread apart, and this compression pulse travels along the spring while each coil moves \emph{back and forth along the axis}. That is a longitudinal wave.

The crucial point for what follows is this: for a transverse wave, the direction of vibration is \emph{not} uniquely fixed by the direction of propagation. If the wave travels horizontally, the vibration may be vertical, horizontal (perpendicular to travel), or in \emph{any} direction lying in the plane perpendicular to propagation. There is a whole family of possible vibration directions. For a longitudinal wave, by contrast, there is no freedom at all: the vibration is forced to lie along the direction of propagation itself. This asymmetry is the deep reason why polarization exists for transverse waves only.

\begin{aretenir}[The key distinction]
For a transverse wave propagating along a given direction, the vibration can take place in \emph{any direction perpendicular} to propagation: infinitely many choices. For a longitudinal wave, the vibration direction is unique: it \emph{is} the direction of propagation.
\end{aretenir}

\courssub{Unpolarized and Plane Polarized Transverse Waves}

Imagine you shake the end of a long rope randomly --- sometimes up and down, sometimes left and right, sometimes diagonally, changing direction unpredictably many times per second. A transverse wave still travels along the rope, but at any fixed point of the rope the vibration direction changes continually and randomly. The vibrations occur in \emph{all planes} containing the direction of propagation. Such a wave is said to be \cle{unpolarized}.

Now imagine instead that you shake the end of the rope strictly up and down, always in the same vertical line. Every particle of the rope then vibrates vertically, and the whole wave --- all its crests and troughs --- lies in one single vertical plane. The vibration has been \emph{restricted to a single plane}. This is the phenomenon of polarization.

\begin{definition}[Polarization of a transverse wave]
A transverse wave is said to be \cle{plane polarized} (or linearly polarized) when the vibrations of the particles of the medium are restricted to a \textbf{single fixed direction} (and therefore to a single plane containing the direction of propagation). Polarization is the operation or phenomenon by which the vibrations of a transverse wave, initially distributed in many directions, are confined to one direction only.

A wave whose vibrations occur randomly in all directions perpendicular to propagation is said to be \cle{unpolarized}.
\end{definition}

\begin{definition}[Plane of polarization]
The \cle{plane of polarization} of a plane polarized wave is the plane that contains both:
\begin{itemize}
\item the \textbf{direction of vibration} of the particles of the medium, and
\item the \textbf{direction of propagation} of the wave.
\end{itemize}
For a rope wave travelling horizontally with vertical vibrations, the plane of polarization is the vertical plane containing the rope.
\end{definition}

\begin{attention}[Do not confuse the two directions]
Students frequently mix up the \emph{direction of vibration} and the \emph{direction of propagation}. They are \textbf{perpendicular} to each other for a transverse wave. The plane of polarization is the plane spanned by these two perpendicular directions --- it is \emph{not} perpendicular to the propagation; on the contrary, it \emph{contains} the direction of propagation.
\end{attention}

\begin{popfigure}
\begin{tikzpicture}[scale=0.9]
% ===== Unpolarized wave (left) =====
\begin{scope}
% Propagation axis
\draw[->, thick, popink] (0,0) -- (5.2,0) node[below left=2pt and 0pt] {\small propagation};
% Representative sinusoid (envelope)
\draw[popblue, thick, domain=0:4.8, samples=120] plot (\x, {0.9*sin(360*\x/1.5)});
\draw[popblue, thick, domain=0:4.8, samples=120] plot (\x, {-0.9*sin(360*\x/1.5)});
% Vibration arrows at a specific point (x=2.4)
\draw[<->, popdark, thick] (2.4,-1.1) -- (2.4,1.1) node[midway, right=2pt] {\small vertical};
\draw[<->, popdark, thick] (2.0,-0.78) -- (2.8,0.78) node[midway, above right=2pt] {\small diagonal};
\draw[<->, popdark, thick] (2.8,-0.78) -- (2.0,0.78) node[midway, above left=2pt] {\small diagonal};
\draw[<->, popdark, thick] (1.95,0) -- (2.85,0) node[midway, above=2pt] {\small horizontal};
\node[align=center, text width=4.8cm] at (2.4,-1.9) {\small \textbf{Unpolarized wave:} vibrations in all planes containing the axis};
\end{scope}
% ===== Polarized wave (right) =====
\begin{scope}[xshift=8.5cm]
\draw[->, thick, popink] (0,0) -- (5.2,0) node[below left=2pt and 0pt] {\small propagation};
% Plane of polarization shaded
\fill[popblueL, opacity=0.35] (0,-1.1) rectangle (4.8,1.1);
% Sinusoid in the plane
\draw[popblue, very thick, domain=0:4.8, samples=120] plot (\x, {0.9*sin(360*\x/1.5)});
\draw[popblue, very thick, domain=0:4.8, samples=120] plot (\x, {-0.9*sin(360*\x/1.5)});
% Single vibration direction
\draw[<->, popdark, thick] (2.4,-1.1) -- (2.4,1.1);
\node[popdark, align=left] at (5.0,0.8) {\small plane of};
\node[popdark, align=left] at (5.0,0.45) {\small polarization};
\node[align=center, text width=4.8cm] at (2.4,-1.9) {\small \textbf{Plane polarized wave:} vibrations confined to one plane};
\end{scope}
\end{tikzpicture}
\legende{An unpolarized transverse wave (left) vibrates in many planes containing the direction of propagation; a plane polarized wave (right) vibrates in one single plane.}
\end{popfigure}

\courssub{Mechanical Demonstration: The Rope and the Two Slits}

The beauty of polarization is that it can be demonstrated with the simplest apparatus imaginable: a rope and two rigid screens, each pierced with a straight slit.

\begin{experience}[The rope-and-slit demonstration]
\textbf{Apparatus:} a long flexible rope stretched horizontally; two rigid screens $S_1$ and $S_2$, each with a straight slit that can be oriented vertically or horizontally; the rope passes freely through both slits.

\textbf{Step 1 --- Parallel slits.} Shake the free end of the rope randomly so that an \emph{unpolarized} transverse wave travels towards $S_1$. Orient the slit of $S_1$ \emph{vertically}.
\begin{itemize}
\item At $S_1$, only the \emph{vertical} component of the vibration can pass through the vertical slit; horizontal vibrations are stopped by the board. Beyond $S_1$, the rope vibrates \emph{vertically only}: the wave is now \textbf{plane polarized} in the vertical plane. $S_1$ acts as a \cle{polarizer}.
\item Place $S_2$ further along the rope with its slit also \emph{vertical}. The vertically polarized wave passes through $S_2$ almost undisturbed: the wave is \textbf{transmitted}. $S_2$ acts as an \cle{analyser}: it tests the state of polarization.
\end{itemize}

\textbf{Step 2 --- Crossed slits.} Keep $S_1$ vertical but rotate the slit of $S_2$ to the \emph{horizontal} position.
\begin{itemize}
\item The wave arriving at $S_2$ vibrates \emph{vertically}, but the slit only allows \emph{horizontal} motion. The board blocks the vibration completely: beyond $S_2$ the rope remains at rest. The wave is \textbf{stopped} --- no transmission.
\end{itemize}

\textbf{Conclusion:} a transverse wave passes through two parallel slits but is blocked by two crossed slits. The first slit polarizes the wave; the second analyses it.
\end{experience}

\begin{popfigure}
\begin{tikzpicture}[scale=0.85]
% ============ TOP: parallel slits ============
\begin{scope}
% Rope before S1 (unpolarized: fuzzy)
\draw[popmuted, thick, domain=0:2.0, samples=100] plot (\x, {0.5*sin(600*\x)});
\draw[popmuted, thick, domain=0:2.0, samples=100] plot (\x, {-0.5*sin(600*\x)});
\draw[popmuted, thick, domain=0:2.0, samples=100] plot (\x, {0.3*sin(600*\x)+0.2});
\draw[popmuted, thick, domain=0:2.0, samples=100] plot (\x, {-0.3*sin(600*\x)-0.2});
% Screen S1 vertical slit
\fill[popink] (2.2,-1.4) rectangle (2.5,-0.3);
\fill[popink] (2.2,0.3) rectangle (2.5,1.4);
\node[above, popink] at (2.35,1.4) {\small $S_1$ (vertical)};
% Polarized wave between S1 and S2 (clean vertical sinusoid)
\draw[popblue, very thick, domain=2.5:4.8, samples=100] plot (\x, {0.6*sin(600*(\x-2.5))});
\draw[popblue, very thick, domain=2.5:4.8, samples=100] plot (\x, {-0.6*sin(600*(\x-2.5))});
% Screen S2 vertical slit
\fill[popink] (5.0,-1.4) rectangle (5.3,-0.3);
\fill[popink] (5.0,0.3) rectangle (5.3,1.4);
\node[above, popink] at (5.15,1.4) {\small $S_2$ (vertical)};
% Transmitted wave
\draw[popblue, very thick, domain=5.3:7.5, samples=100] plot (\x, {0.6*sin(600*(\x-5.3))});
\draw[popblue, very thick, domain=5.3:7.5, samples=100] plot (\x, {-0.6*sin(600*(\x-5.3))});
\draw[->, popdark, thick] (7.5,0) -- (8.2,0);
\node[popgreen, align=center] at (3.65,-1.8) {\small \textbf{Parallel slits:} wave transmitted};
\end{scope}
% ============ BOTTOM: crossed slits ============
\begin{scope}[yshift=-5.0cm]
% Rope before S1 (unpolarized)
\draw[popmuted, thick, domain=0:2.0, samples=100] plot (\x, {0.5*sin(600*\x)});
\draw[popmuted, thick, domain=0:2.0, samples=100] plot (\x, {-0.5*sin(600*\x)});
\draw[popmuted, thick, domain=0:2.0, samples=100] plot (\x, {0.3*sin(600*\x)+0.2});
\draw[popmuted, thick, domain=0:2.0, samples=100] plot (\x, {-0.3*sin(600*\x)-0.2});
% Screen S1 vertical slit
\fill[popink] (2.2,-1.4) rectangle (2.5,-0.3);
\fill[popink] (2.2,0.3) rectangle (2.5,1.4);
\node[above, popink] at (2.35,1.4) {\small $S_1$ (vertical)};
% Polarized wave between S1 and S2
\draw[popblue, very thick, domain=2.5:4.8, samples=100] plot (\x, {0.6*sin(600*(\x-2.5))});
\draw[popblue, very thick, domain=2.5:4.8, samples=100] plot (\x, {-0.6*sin(600*(\x-2.5))});
% Screen S2 horizontal slit (gap in middle horizontally)
% Left block
\fill[popink] (5.0,-1.4) rectangle (4.7,1.4);
% Right block
\fill[popink] (5.6,-1.4) rectangle (5.3,1.4);
% Top block
\fill[popink] (5.0,0.3) rectangle (5.3,1.4);
% Bottom block
\fill[popink] (5.0,-1.4) rectangle (5.3,-0.3);
\node[above, popink] at (5.15,1.4) {\small $S_2$ (horizontal)};
% No wave after (flat line)
\draw[popblue, thick] (5.3,0) -- (8.2,0);
\node[poporange, align=center] at (3.65,-1.8) {\small \textbf{Crossed slits:} wave blocked};
\node[popmuted, font=\small] at (6.8,0.4) {no vibration};
\end{scope}
\end{tikzpicture}
\legende{The rope-and-slit model. Top: two parallel vertical slits --- the vertically polarized wave is transmitted. Bottom: crossed slits (vertical then horizontal) --- the horizontal slit blocks the vertical vibration and nothing passes.}
\end{popfigure}

The vocabulary introduced here is used throughout physics and at the GCE examination:

\begin{definition}[Polarizer and analyser]
A \cle{polarizer} is a device which, receiving an unpolarized transverse wave, transmits only the component of vibration parallel to one privileged direction: it produces a plane polarized wave. In the rope model, the polarizer is the first slit.

An \cle{analyser} is a device, identical in nature to the polarizer, used to \emph{detect} the state of polarization of the wave: rotating it and observing whether the wave is transmitted or blocked reveals the plane of polarization. In the rope model, the analyser is the second slit.
\end{definition}

\begin{methode}[Predicting transmission with the rope-and-slit model]
To decide whether a transverse wave is transmitted through a system of slits (or polarizers), proceed step by step:
\begin{enumerate}
\item \textbf{Identify the state of the incident wave.} Is it unpolarized (vibrations in all directions) or already plane polarized? If polarized, note the direction of vibration.
\item \textbf{Apply the first slit (polarizer).} Only the component of vibration \emph{parallel to the slit} passes. After the slit, the wave is plane polarized along the slit direction.
\item \textbf{Apply the second slit (analyser).} Compare the vibration direction of the wave arriving at the analyser with the orientation of its slit.
\item \textbf{Conclude:}
\begin{itemize}
\item slit \emph{parallel} to the vibration $\Rightarrow$ wave \textbf{transmitted} (maximum);
\item slit \emph{perpendicular} to the vibration (crossed) $\Rightarrow$ wave \textbf{completely blocked};
\item slit at an intermediate angle $\theta$ $\Rightarrow$ only the \emph{component} of the vibration along the slit is transmitted (amplitude reduced by factor $\cos\theta$, intensity by $\cos^2\theta$).
\end{itemize}
\end{enumerate}
\end{methode}

\begin{exoresolu}[Two slits on a rope]
A long horizontal rope carries a transverse wave. The rope passes through slit $A$, which is vertical, and then through slit $B$, which makes an angle of $90^\circ$ with slit $A$ (i.e.\ it is horizontal). The free end of the rope is shaken randomly.

\textbf{1.} What is the state of polarization of the wave between $A$ and $B$?

\textbf{2.} Is any wave observed beyond $B$? Justify.

\textbf{3.} Slit $B$ is now rotated until it is vertical. What happens?
\tcblower
\textbf{1.} The incident wave is unpolarized: the random shaking produces vibrations in all directions perpendicular to the rope. Slit $A$ is vertical, so it transmits only the \emph{vertical} component of the vibration. Between $A$ and $B$ the rope vibrates vertically only: the wave is \textbf{plane polarized in the vertical plane}, with $A$ acting as polarizer.

\textbf{2.} Slit $B$ is horizontal, i.e.\ \emph{crossed} with the vertical vibration arriving at it. A vertical vibration cannot pass through a horizontal slit, so the wave is \textbf{completely blocked}: the rope beyond $B$ remains at rest. No wave is observed beyond $B$.

\textbf{3.} With $B$ vertical, its slit is now \emph{parallel} to the vertical vibration of the polarized wave. The wave is \textbf{transmitted} through $B$ with essentially its full amplitude: a vertical vibration is observed on the rope beyond $B$. This is exactly the crossed-versus-parallel behaviour predicted by the rope-and-slit model.
\end{exoresolu}

\courssub{Why Longitudinal Waves Cannot Be Polarized}

Let us now ask the question the other way round: could a sound wave in air be polarized? Suppose we try to build the analogue of the rope-and-slit experiment for sound. Sound in air is a longitudinal wave: the air molecules oscillate \emph{back and forth along the direction of propagation}. Whatever ``filter'' we imagine placing across the beam, the vibration direction is already fixed --- it is the direction of travel itself. There is no transverse component to select, no orientation to restrict, and rotating any analyser around the direction of propagation changes nothing, because the situation is perfectly symmetric about that axis.

\begin{propriete}[Polarization is the test of transversality]
Only \cle{transverse waves} can be polarized. A longitudinal wave (such as sound in air) cannot be polarized, because its vibration direction is uniquely fixed along the direction of propagation: there is no freedom of orientation to restrict.

Conversely, if a wave phenomenon can be polarized --- if an analyser can be found which, when rotated, extinguishes the wave --- then the wave is necessarily \textbf{transverse}. Polarization is therefore the decisive \emph{experimental proof of the transverse nature of a wave}. (This reasoning is examinable: state it clearly and justify it with the rope-and-slit model.)
\end{propriete}

This is precisely how physicists established that light is a transverse wave. Light cannot be seen to vibrate directly, but it can be polarized --- as the Douala rider's sunglasses prove --- and polarization is only possible for transverse vibrations. Sound in air, which can never be polarized, is confirmed to be longitudinal.

\begin{attention}[Common mistake: ``polarizing'' sound]
A very frequent error at the GCE is to claim that sound waves can be polarized, or to apply the polarizer--analyser description to sound in air. This is \textbf{wrong}: sound in air (and in fluids generally) is \emph{longitudinal}; the vibration is along the direction of propagation, so no slit or filter can extinguish it by rotation. Never write that ``the polarizer blocks the sound''. The correct statement is: \emph{sound in air cannot be polarized because it is a longitudinal wave.}
\end{attention}

\begin{savaistu}[Polarized sunglasses and glare]
Light reflected by a horizontal surface --- a wet road in Douala, the bonnet or windscreen of a car, the surface of the Wouri river --- becomes partially polarized, with its vibration preferentially \emph{horizontal}. Polarized sunglasses contain a plastic film whose molecules act like a microscopic ``slit'' oriented \emph{vertically}: they transmit the vertical component and absorb the horizontal glare. Tilt your head (rotate the ``analyser'') and the glare returns --- the everyday version of the crossed-slits experiment! Ordinary tinted glasses merely dim everything without selecting any vibration direction, which is why the rider's friend keeps squinting.
\end{savaistu}

\begin{exercice}[Testing your understanding]
\textbf{1.} Define a plane polarized transverse wave and state what is meant by its plane of polarization.

\textbf{2.} In the rope-and-slit experiment, the first slit is vertical and the second makes an angle of $90^\circ$ with the first. Explain, step by step, why no vibration is observed beyond the second slit.

\textbf{3.} A student claims: ``Since sound can be blocked by a wall, it can also be blocked by an analyser, so sound can be polarized.'' Identify the two errors in this statement and correct them.

\textbf{4.} Explain why the fact that light can be polarized proves that light is a transverse wave.
\end{exercice}

\begin{corrige}[Exercise 1]
\textbf{1.} A plane polarized transverse wave is a wave whose particle vibrations are restricted to a single fixed direction perpendicular to the propagation. The plane of polarization is the plane containing the direction of vibration and the direction of propagation.

\textbf{2.} The randomly shaken rope sends an unpolarized wave onto the first (vertical) slit, which transmits only the vertical component: the wave becomes plane polarized vertically. Arriving at the second (horizontal) slit, the vertical vibration is perpendicular to the slit and cannot pass; the board absorbs it. The rope beyond the second slit therefore does not vibrate.

\textbf{3.} First error: blocking a wave by absorption (a wall) is not polarization; polarization is a \emph{selection of a vibration direction}, not a mere stopping of energy. Second error: sound in air is longitudinal, its vibration lies along the propagation direction, so rotating any analyser cannot extinguish it. Sound in air cannot be polarized.

\textbf{4.} Only transverse waves possess a vibration direction that can vary around the propagation axis and can therefore be restricted to one plane. Since experiments (polarizer and analyser, extinction by crossed polarizers) succeed with light, light must have a transverse vibration: light is a transverse wave.
\end{corrige}

\begin{aretenir}[Summary of Section 1]
\begin{itemize}
\item A transverse wave vibrates \emph{perpendicular} to propagation; a longitudinal wave vibrates \emph{parallel} to propagation.
\item \cle{Polarization} is the restriction of the vibrations of a transverse wave to a single direction; the \cle{plane of polarization} contains the vibration direction and the propagation direction.
\item An unpolarized wave vibrates in all planes containing the direction of propagation.
\item The rope-and-slit experiment models the \cle{polarizer} (first slit) and the \cle{analyser} (second slit): parallel slits transmit, crossed slits block.
\item Only transverse waves can be polarized; polarization is the experimental \emph{test of transversality}. Sound in air, being longitudinal, can never be polarized.
\end{itemize}
\end{aretenir}

\coursec{Production and Detection of Plane Polarized Light; Applications}

In the previous section we established, with the help of a rope and a pair of slits, that a transverse wave can be \cle{polarized}: its vibrations can be confined to a single plane. We also saw that a longitudinal wave (sound, for example) can never be polarized. We now turn to light. The fact that light \emph{can} be polarized was historically one of the strongest proofs that light is a \cle{transverse wave}. In this section we explain what polarized light is, how it is produced, how it is detected, and why the phenomenon is so useful in everyday technology.

\courssub{Light as a Transverse Electromagnetic Wave}

Light is an electromagnetic wave: it consists of an electric field $\vec{E}$ and a magnetic field $\vec{B}$ which vibrate \emph{perpendicular to each other} and \emph{perpendicular to the direction of propagation}. No material medium is needed: light travels through the vacuum of space from the Sun to Cameroon every day.

Because the vibrations are transverse, they can be polarized. By convention, the \cle{plane of polarization} of a light wave is defined by the plane containing the \cle{electric field vector} $\vec{E}$ and the direction of propagation. If $\vec{E}$ always vibrates along one fixed direction --- say vertically --- while the wave travels horizontally, the light is said to be \cle{plane polarized} (or linearly polarized) in the vertical plane.

\begin{definition}[Unpolarized (natural) light]
Ordinary light --- from the Sun, a candle, or an electric bulb --- is \cle{unpolarized}. It is emitted by a huge number of atoms vibrating independently, so the electric field $\vec{E}$ vibrates in \textbf{all possible directions} perpendicular to the direction of propagation, changing randomly millions of times per second. There is no preferred plane of vibration.
\end{definition}

\begin{popfigure}
\begin{center}
\begin{tikzpicture}[scale=1.0]
% Unpolarized light: several E directions
\draw[->,thick,popink] (0,0) -- (4.2,0) node[below left,font=\small]{propagation};
\foreach \a in {20,65,110,155}{
  \draw[->,popblue,thick] (1.2,0) -- ++(\a:0.75);
  \draw[->,popblue,thick] (1.2,0) -- ++(\a+180:0.75);
}
\node[popblue,font=\small,align=center,text width=2.6cm] at (1.2,-1.35) {unpolarized: $\vec{E}$ in all transverse directions};
% Polarized light: one direction
\draw[->,thick,popink] (5.6,0) -- (9.8,0) node[below left,font=\small]{propagation};
\draw[->,poporange,very thick] (7.0,0) -- (7.0,0.85);
\draw[->,poporange,very thick] (7.0,0) -- (7.0,-0.85);
\draw[->,poporange,very thick] (8.2,0) -- (8.2,0.85);
\draw[->,poporange,very thick] (8.2,0) -- (8.2,-0.85);
\node[poporange,font=\small,align=center,text width=2.6cm] at (7.6,-1.35) {plane polarized: $\vec{E}$ along one direction};
\end{tikzpicture}
\legende{Unpolarized light (left): $\vec{E}$ vibrates in every transverse direction. Plane polarized light (right): $\vec{E}$ vibrates along a single direction.}
\end{center}
\end{popfigure}

\courssub{Producing Plane Polarized Light: the Polaroid Filter}

The simplest way to polarize light is to pass it through a \cle{polaroid} sheet. A polaroid is a plastic film containing long-chain molecules all aligned in the same direction (a process called \emph{dichroism}). These molecules strongly absorb the component of $\vec{E}$ that vibrates \emph{parallel} to the chains, while transmitting the component vibrating \emph{perpendicular} to them. The direction of the transmitted vibration is called the \cle{transmission axis} of the polaroid.

\begin{definition}[Polarizer and analyser]
A \cle{polarizer} is a device (usually a polaroid sheet) that transforms unpolarized light into plane polarized light, by absorbing all vibrations except those parallel to its transmission axis.
An \cle{analyser} is a second polaroid, placed after the polarizer, used to \emph{detect} and \emph{measure} the state of polarization of the light. Rotating the analyser changes the transmitted intensity, which reveals the plane of polarization.
\end{definition}

\begin{attention}[The polarizer does not ``create'' the vibration]
The polarizer does not force light to vibrate in a new direction. It \emph{selects}, from the randomly oriented vibrations of natural light, the component parallel to its transmission axis and absorbs the rest. Consequently, an ideal polarizer transmits only \textbf{half} of the incident unpolarized intensity: $I_0 = \tfrac{1}{2}I_{\text{incident}}$.
\end{attention}

\textbf{Other production methods (GCE-useful extras).}
\begin{itemize}
\item \emph{Polarization by reflection.} When unpolarized light reflects from a non-metallic surface (water, glass, a road), the reflected light is partially polarized parallel to the surface. At a particular angle of incidence, the \cle{Brewster angle} $\theta_B$, the reflected light is \emph{completely} plane polarized and the reflected and refracted rays are perpendicular to each other. The Brewster angle satisfies $\tan\theta_B = n_2/n_1$, where $n_1$ and $n_2$ are the refractive indices of the two media. This is why glare from a wet road near Douala on a sunny afternoon is strongly polarized.
\item \emph{Polarization by scattering.} Blue skylight, scattered by air molecules, is partially polarized. Looking at the sky at $90^{\circ}$ from the Sun through a polaroid and rotating it, the brightness changes noticeably.
\end{itemize}

\courssub{Detecting Polarized Light: the Analyser and Crossed Polaroids}

Place a second polaroid (the analyser) after the polarizer and rotate it slowly about the direction of propagation:
\begin{itemize}
\item when the transmission axes are \textbf{parallel} ($\theta = 0^{\circ}$), the transmitted intensity is \textbf{maximum};
\item as the analyser rotates, the intensity decreases continuously;
\item when the axes are \textbf{perpendicular} ($\theta = 90^{\circ}$), the transmitted intensity is \textbf{zero}: the polaroids are said to be \emph{crossed} and no light passes through.
\end{itemize}

\begin{propriete}[Crossed polaroids]
Two polaroids whose transmission axes are perpendicular ($\theta = 90^{\circ}$) transmit \textbf{no light}: the polarizer passes only vertical vibrations, and the analyser passes only horizontal ones, so nothing survives. This complete extinction is the experimental signature that light is a transverse wave.
\end{propriete}

\begin{popfigure}
\begin{center}
\begin{tikzpicture}[scale=0.95]
% incoming unpolarized light
\foreach \a in {30,90,150}{
  \draw[->,popblue,thick] (0.4,0) -- ++(\a:0.6);
  \draw[->,popblue,thick] (0.4,0) -- ++(\a+180:0.6);
}
\draw[->,thick,popink] (-0.4,0) -- (1.6,0);
\node[font=\small,popblue] at (0.4,1.05) {unpolarized};
% Polarizer
\draw[fill=popblueL,draw=popblue,thick] (2.0,-1.1) rectangle (2.35,1.1);
\foreach \y in {-0.9,-0.45,0,0.45,0.9}{\draw[popblue] (2.17,\y-0.13) -- (2.17,\y+0.13);}
\node[font=\small,align=center] at (2.17,-1.5) {polarizer\\ (vertical axis)};
% polarized light between
\draw[->,thick,popink] (2.35,0) -- (4.6,0);
\draw[->,poporange,very thick] (3.0,0) -- (3.0,0.7);
\draw[->,poporange,very thick] (3.0,0) -- (3.0,-0.7);
\draw[->,poporange,very thick] (3.9,0) -- (3.9,0.7);
\draw[->,poporange,very thick] (3.9,0) -- (3.9,-0.7);
\node[font=\small,poporange] at (3.5,1.05) {plane polarized};
% Analyser crossed
\draw[fill=popgreenL,draw=popgreen,thick] (4.8,-1.1) rectangle (5.15,1.1);
\foreach \x in {-0.9,-0.45,0,0.45,0.9}{\draw[popgreen] (4.97-0.13,\x) -- (4.97+0.13,\x);}
\node[font=\small,align=center] at (4.97,-1.5) {analyser\\ (horizontal axis)};
% no light after
\draw[->,thick,popink,dashed] (5.15,0) -- (6.3,0);
\node[font=\small,popdark] at (6.9,0) {no light};
\draw[popdark,thick] (5.6,0.35) -- (5.9,-0.35);
\draw[popdark,thick] (5.6,-0.35) -- (5.9,0.35);
\end{tikzpicture}
\legende{Crossed polaroids: the polarizer produces vertically polarized light; the analyser, with a horizontal transmission axis, absorbs it completely. The transmitted intensity falls to zero.}
\end{center}
\end{popfigure}

\courssub{Malus's Law (Examinable)}

Let plane polarized light of intensity $I_0$ fall on an analyser whose transmission axis makes an angle $\theta$ with the plane of polarization of the incident light. Only the component of the electric field \emph{parallel} to the analyser's axis is transmitted:
\[ E_{\text{transmitted}} = E_0\cos\theta . \]
Since the intensity of an electromagnetic wave is proportional to the \emph{square} of the amplitude of its electric field, $I \propto E^2$, we obtain:

\begin{propriete}[Malus's law]
The intensity transmitted by an analyser whose axis makes an angle $\theta$ with the plane of polarization of the incident polarized light is
\[ I = I_0\cos^2\theta \]
where $I_0$ is the intensity incident on the analyser. In particular $I=I_0$ for $\theta=0^{\circ}$ (parallel axes) and $I=0$ for $\theta=90^{\circ}$ (crossed axes).
\end{propriete}

\begin{popfigure}
\begin{center}
\begin{tikzpicture}
\begin{axis}[
  width=11cm,height=6.2cm,
  xlabel={$\theta$ (degrees)}, ylabel={$I/I_0$},
  xmin=0,xmax=180,ymin=0,ymax=1.1,
  xtick={0,30,60,90,120,150,180},
  ytick={0,0.25,0.5,0.75,1},
  grid=both, grid style={popline!40},
  axis line style={popink},
  label style={popink},
]
\addplot[domain=0:180,samples=120,very thick,poporange] {cos(x)^2};
\addplot[mark=*,only marks,popblue] coordinates {(0,1)(90,0)(180,1)};
\node[font=\small,popblue] at (axis cs:20,1.06) {parallel};
\node[font=\small,popblue] at (axis cs:90,0.09) {crossed};
\end{axis}
\end{tikzpicture}
\legende{Malus's law: transmitted intensity $I=I_0\cos^2\theta$ as a function of the angle $\theta$ between the transmission axes of polarizer and analyser.}
\end{center}
\end{popfigure}

\begin{methode}[Identifying the polarizer and the analyser]
In any two-polaroid problem:
\begin{enumerate}
\item The \textbf{first} filter met by the (unpolarized) light is the \cle{polarizer}. It halves the intensity: $I_0 = \tfrac{1}{2}I_{\text{incident}}$.
\item The \textbf{second} filter is the \cle{analyser}. Apply Malus's law between the two: $I = I_0\cos^2\theta$.
\item The angle $\theta$ is measured \textbf{between the two transmission axes}, never from the direction of propagation.
\item If a third polaroid is inserted, apply Malus's law step by step between consecutive filters.
\end{enumerate}
\end{methode}

\begin{exoresolu}[Two crossed polaroids with a third inserted]
Unpolarized light of intensity $I_i = 8.0\ \mathrm{W\,m^{-2}}$ falls on a polarizer $P_1$. A second polaroid $P_2$ is crossed with $P_1$ ($90^{\circ}$). (a) What intensity emerges from $P_2$? (b) A third polaroid $P_3$ is inserted between them, with its axis at $45^{\circ}$ to that of $P_1$. Find the final transmitted intensity.
\tcblower
\textbf{(a)} With $P_1$ and $P_2$ crossed, Malus's law gives
\[ I = I_0\cos^2 90^{\circ} = 0. \]
No light emerges: $I = 0\ \mathrm{W\,m^{-2}}$.

\textbf{(b)} Step by step:
\begin{align*}
\text{After } P_1:\quad & I_1 = \tfrac{1}{2}I_i = \tfrac{1}{2}(8.0) = 4.0\ \mathrm{W\,m^{-2}}\\
\text{After } P_3\ (\theta=45^{\circ}):\quad & I_2 = I_1\cos^2 45^{\circ} = 4.0\times(0.7071)^2 = 2.0\ \mathrm{W\,m^{-2}}\\
\text{After } P_2\ (\theta=45^{\circ}\text{ to }P_3):\quad & I_3 = I_2\cos^2 45^{\circ} = 2.0\times 0.5 = 1.0\ \mathrm{W\,m^{-2}}
\end{align*}
\emph{Remark:} inserting $P_3$ at $45^{\circ}$ makes light pass through a previously opaque system --- a striking confirmation of the vector nature of $\vec{E}$.
\end{exoresolu}

\begin{attention}[Do not confuse the angles]
The angle $\theta$ in Malus's law is the angle \textbf{between the transmission axes} of the polarizer and the analyser (or between the plane of polarization and the analyser axis). It has \emph{nothing to do} with the angle of incidence of the light on the filter. Also remember: Malus's law applies to light that is \emph{already polarized}; for unpolarized light falling on the first polaroid, simply divide the intensity by two.
\end{attention}

\courssub{Applications of Polarized Light}

\begin{itemize}
\item \textbf{Polarized sunglasses.} Light reflected by horizontal surfaces (a wet road, the surface of the Wouri river) is partially polarized horizontally. The polaroid lenses of the glasses have a \emph{vertical} transmission axis, so they block this glare while transmitting a good fraction of the useful light. Fishermen and drivers benefit greatly.
\item \textbf{LCD screens.} A liquid-crystal display consists of a layer of liquid crystal sandwiched between two crossed polaroids. Applying a voltage rotates the plane of polarization locally, controlling pixel by pixel whether light passes --- the basis of calculators, phones and television screens.
\item \textbf{Photoelastic stress analysis.} Certain transparent plastics become doubly refracting when stressed. Placed between crossed polaroids, a model of a bridge or gear shows coloured fringes that reveal the zones of greatest mechanical stress --- a tool used by engineers before construction.
\item \textbf{3D cinema.} The two images (left eye and right eye) are projected with perpendicular planes of polarization; the viewer wears glasses whose two lenses are crossed polaroids, so each eye sees only its own image, creating the impression of depth.
\end{itemize}

\begin{savaistu}[Polarization in daily life]
If you tilt your head while wearing polarized sunglasses and looking at an LCD screen, the screen appears to go dark at a certain angle: the screen's light is already polarized, and you have effectively ``crossed'' your glasses with it. This is a quick home test of both polarization and Malus's law!
\end{savaistu}

\begin{aretenir}[Key points]
\begin{itemize}
\item Light is a transverse electromagnetic wave; the direction of $\vec{E}$ defines the plane of polarization.
\item Natural light is unpolarized: $\vec{E}$ vibrates in all transverse directions.
\item A polarizer produces plane polarized light by selective absorption and transmits half the incident unpolarized intensity.
\item An analyser detects polarization: intensity maximum for parallel axes, zero for crossed axes ($90^{\circ}$).
\item Malus's law: $I = I_0\cos^2\theta$, with $\theta$ the angle between transmission axes.
\item Applications: sunglasses, LCD screens, stress analysis, 3D cinema.
\end{itemize}
\end{aretenir}

\begin{exercice}[Polaroids at 30 degrees]
Unpolarized light of intensity $12\ \mathrm{W\,m^{-2}}$ passes through a polarizer and then an analyser whose axis is at $30^{\circ}$ to that of the polarizer. Calculate the intensity transmitted by the analyser.
\end{exercice}

\begin{corrige}[Polaroids at 30 degrees]
After the polarizer: $I_0 = \tfrac{1}{2}(12) = 6.0\ \mathrm{W\,m^{-2}}$.
Malus's law: $I = I_0\cos^2 30^{\circ} = 6.0\times(0.866)^2 = 6.0\times 0.75 = 4.5\ \mathrm{W\,m^{-2}}$.
\end{corrige}

\coursec{Speed of Transverse Waves on a Taut String: Factors and Formula}

In the previous section we studied how transverse waves can be polarised, and we insisted that the speed of a wave is a property of the \emph{medium} that carries it. We now return to mechanical waves and ask a very concrete question: \emph{what determines the speed of a transverse wave travelling along a stretched string?} This question is not academic. Every guitarist in Douala, every piano tuner, and every engineer who designs suspension cables must know the answer. In this section we identify the factors experimentally, derive the form of the formula by dimensional analysis, and learn to use it confidently in numerical problems.

\courssub{Qualitative Observations: What Does the Speed Depend On?}

Take a guitar string and pluck it gently. A transverse disturbance runs along the string, is reflected at the fixed ends, and sets the string vibrating. Now perform two simple thought experiments (or, better, real ones in the school laboratory with a long rubber cord).

\textbf{Observation 1: effect of tension.} Tighten the string by turning the tuning peg, so that the tension $T$ increases. Pluck it again: the note becomes higher (the frequency increases) and, if you watch a pulse travel along a long rope, you see that the disturbance now travels \emph{faster}. Slacken the string and the pulse becomes sluggish. A loose rope can barely transmit a transverse wave at all, because there is almost no restoring force pulling a displaced element back towards the equilibrium position.

\textbf{Observation 2: effect of the mass of the string.} Compare a thin nylon string with a thick steel cable under the same tension. The pulse travels noticeably \emph{slower} along the heavy cable. The heavier the string (per unit length), the more inertia each element of the string has, and the harder it is for the tension to accelerate it up and down.

These observations point to exactly two factors:
\begin{itemize}
\item the \cle{tension} $T$ in the string (in newtons): increasing $T$ increases the wave speed;
\item the \cle{mass per unit length} of the string, which measures its inertia: increasing it decreases the wave speed.
\end{itemize}

\begin{definition}[Linear density (mass per unit length)]
The \cle{linear density} of a uniform string or wire is its mass divided by its length:
\[
\mu = \frac{m}{L}
\]
where $m$ is the mass of the string (in kg) and $L$ its length (in m). The SI unit of $\mu$ is the kilogram per metre, $\mathrm{kg\,m^{-1}}$.
\end{definition}

\begin{exemplebox}[Linear density of a guitar string]
A steel guitar string has length $L = \SI{0.65}{m}$ and mass $m = \SI{3.9e-3}{kg}$. Its linear density is
\[
\mu = \frac{m}{L} = \frac{3.9\times 10^{-3}}{0.65} = \SI{6.0e-3}{kg.m^{-1}}.
\]
A bass string of the same length but mass $\SI{1.3e-2}{kg}$ has $\mu = \SI{2.0e-2}{kg.m^{-1}}$: it is more than three times ``heavier per metre'', which is why it produces lower notes under comparable tension.
\end{exemplebox}

\courssub{Dimensional Analysis: Finding the Form of the Formula}

We suspect that the wave speed $v$ depends only on $T$ and $\mu$. Let us suppose a relation of the form
\[
v = k\,T^{a}\,\mu^{b},
\]
where $k$ is a dimensionless constant and $a$, $b$ are powers to be found. We use the dimensions of each quantity:
\[
[v] = \mathrm{L\,T^{-1}}, \qquad [T] = \mathrm{M\,L\,T^{-2}}, \qquad [\mu] = \mathrm{M\,L^{-1}}.
\]
(Be careful: the symbol $T$ is used for tension, while $\mathrm{T}$ in brackets denotes the dimension of time.) Equating dimensions on both sides:
\[
\mathrm{L\,T^{-1}} = \left(\mathrm{M\,L\,T^{-2}}\right)^{a}\left(\mathrm{M\,L^{-1}}\right)^{b}
= \mathrm{M}^{a+b}\,\mathrm{L}^{a-b}\,\mathrm{T}^{-2a}.
\]
Matching the exponents of each dimension:
\begin{align*}
\text{Mass:} \quad & a + b = 0,\\
\text{Length:} \quad & a - b = 1,\\
\text{Time:} \quad & -2a = -1 \;\Rightarrow\; a = \tfrac{1}{2}.
\end{align*}
From $a = \tfrac{1}{2}$ we get $b = -\tfrac{1}{2}$. Therefore
\[
v = k\,T^{1/2}\,\mu^{-1/2} = k\sqrt{\frac{T}{\mu}}.
\]
Dimensional analysis cannot give the dimensionless constant $k$; a full mechanical treatment (applying Newton's second law to a small element of the string) shows that $k = 1$ exactly. Note the consistency check:
\[
\left[\frac{T}{\mu}\right] = \frac{\mathrm{M\,L\,T^{-2}}}{\mathrm{M\,L^{-1}}} = \mathrm{L^{2}\,T^{-2}},
\qquad \sqrt{\mathrm{L^{2}\,T^{-2}}} = \mathrm{L\,T^{-1}} = [v]. \ \checkmark
\]

\begin{methode}[Dimensional-analysis check before applying any wave-speed formula]
\begin{enumerate}
\item Write the formula you intend to use, e.g. $v = \sqrt{T/\mu}$.
\item Replace each quantity by its dimensions: $[T] = \mathrm{M\,L\,T^{-2}}$, $[\mu] = \mathrm{M\,L^{-1}}$, $[v] = \mathrm{L\,T^{-1}}$.
\item Simplify the right-hand side: $\dfrac{\mathrm{M\,L\,T^{-2}}}{\mathrm{M\,L^{-1}}} = \mathrm{L^{2}\,T^{-2}}$, whose square root is $\mathrm{L\,T^{-1}}$.
\item Compare with the left-hand side. If the dimensions agree, the formula is at least dimensionally correct; if not, the formula is certainly wrong (for example $v = \sqrt{\mu/T}$ fails the test immediately).
\item Repeat the check with SI units in every numerical calculation: $T$ in newtons, $\mu$ in $\mathrm{kg\,m^{-1}}$, giving $v$ in $\mathrm{m\,s^{-1}}$.
\end{enumerate}
\end{methode}

\courssub{The Wave-Speed Formula}

\begin{propriete}[Speed of a transverse wave on a taut string]
The speed of a transverse wave travelling along a taut string of linear density $\mu$ under tension $T$ is
\[
\boxed{\,v = \sqrt{\dfrac{T}{\mu}}\,}
\]
where
\begin{itemize}
\item $v$ is the wave speed, in $\mathrm{m\,s^{-1}}$;
\item $T$ is the tension in the string, in newtons (N);
\item $\mu = m/L$ is the mass per unit length, in $\mathrm{kg\,m^{-1}}$.
\end{itemize}
The derivation from Newton's second law applied to a string element is not required at this level, but the dimensional-analysis derivation above \emph{is} examinable.
\end{propriete}

\begin{popfigure}
\begin{center}
\begin{tikzpicture}[scale=1.0, >=stealth]
% supports
\fill[poppurple] (-0.15,-0.6) rectangle (0.15,0.6);
\fill[poppurple] (9.85,-0.6) rectangle (10.15,0.6);
\node[below, text=popmuted] at (0,-0.65) {fixed end};
\node[below, text=popmuted] at (10,-0.65) {fixed end};
% string
\draw[very thick, popink] (0,0) -- (10,0);
% pulse
\draw[very thick, poporange, domain=3.2:5.2, samples=60] plot (\x, {0.8*exp(-((\x-4.2)^2)/0.18)});
% velocity arrow
\draw[->, very thick, popblue] (5.4,0.85) -- (7.0,0.85) node[midway, above]{$v = \sqrt{T/\mu}$};
% tension arrows
\draw[->, very thick, popgreen] (0.2,0.15) -- (1.6,0.15) node[midway, above]{$T$};
\draw[->, very thick, popgreen] (9.8,0.15) -- (8.4,0.15) node[midway, above]{$T$};
% label
\node[text=popmuted] at (2.0,-0.45) {string of linear density $\mu = m/L$};
\node[text=poporange] at (4.2,1.35) {pulse};
\end{tikzpicture}
\end{center}
\legende{A pulse travelling with speed $v$ along a taut string fixed at both ends. The tension $T$ pulls on every element of the string; the linear density $\mu$ measures its inertia.}
\end{popfigure}

\begin{propriete}[The medium fixes the speed, not the source]
For a \emph{given} string under a \emph{given} tension, the speed $v = \sqrt{T/\mu}$ is fixed once and for all. It does \textbf{not} depend on:
\begin{itemize}
\item the \cle{frequency} $f$ of the wave;
\item the \cle{amplitude} of the disturbance (for small oscillations);
\item the shape of the pulse or the way the string was plucked.
\end{itemize}
If the source vibrates faster, it is the \emph{wavelength} $\lambda = v/f$ that adjusts, not the speed. This is a general feature of waves in non-dispersive media: \textbf{the medium determines the speed}.
\end{propriete}

\courssub{Combining with the Wave Equation $v = f\lambda$}

Every progressive wave also obeys the universal relation $v = f\lambda$. Combining this with $v = \sqrt{T/\mu}$ gives a powerful tool linking the four quantities:
\[
f\lambda = \sqrt{\frac{T}{\mu}}
\qquad\Longleftrightarrow\qquad
f = \frac{1}{\lambda}\sqrt{\frac{T}{\mu}}.
\]
For a string of length $L$ fixed at both ends and vibrating in its fundamental mode, the length of the string is half a wavelength, $\lambda = 2L$, so the fundamental frequency is
\[
f_1 = \frac{1}{2L}\sqrt{\frac{T}{\mu}}.
\]
This single formula explains the whole design of stringed instruments:
\begin{itemize}
\item \textbf{Tightening} a string (increasing $T$) raises the pitch: this is what a musician does with the tuning pegs.
\item \textbf{Shortening} the vibrating part of the string (smaller $L$) gives higher notes: this is why the player presses the string against the frets.
\item \textbf{Thinner} strings (smaller $\mu$) give higher notes: this is why the treble strings of a piano or guitar are thin and the bass strings are thick or wound with extra wire.
\end{itemize}

\begin{popfigure}
\begin{center}
\begin{tikzpicture}
\begin{axis}[
  width=0.72\linewidth, height=6cm,
  axis lines=left,
  xlabel={$\sqrt{T}\;(\mathrm{N^{1/2}})$},
  ylabel={$v\;(\mathrm{m\,s^{-1}})$},
  xlabel style={text=popink}, ylabel style={text=popink},
  xmin=0, xmax=10.5, ymin=0, ymax=105,
  xtick={0,2,4,6,8,10}, ytick={0,20,40,60,80,100},
  grid=both, grid style={popline, very thin},
  tick label style={font=\small},
]
\addplot[very thick, popblue, domain=0:10] {10*x};
\addplot[only marks, mark=*, mark size=2pt, poporange] coordinates {(4,40) (6,60) (8,80) (10,100)};
\node[text=popblue, anchor=south west, font=\small] at (axis cs:4.4,72) {$v = \dfrac{1}{\sqrt{\mu}}\sqrt{T}$};
\node[text=popmuted, font=\small, align=center] at (axis cs:7.2,30) {straight line\\ through origin\\ (fixed $\mu$)};
\end{axis}
\end{tikzpicture}
\end{center}
\legende{For a fixed string (constant $\mu$), the wave speed $v$ is proportional to $\sqrt{T}$: the graph of $v$ against $\sqrt{T}$ is a straight line through the origin, of slope $1/\sqrt{\mu}$. Equivalently, a graph of $v^{2}$ against $T$ is a straight line of slope $1/\mu$.}
\end{popfigure}

\courssub{Worked Examples}

\begin{exoresolu}[Speed of a pulse on a guitar string]
A guitar string has length $L = \SI{0.65}{m}$ and mass $m = \SI{3.9e-3}{kg}$. It is stretched to a tension $T = \SI{96}{N}$. Calculate the speed of a transverse wave on this string.
\tcblower
\textbf{Step 1: linear density.}
\[
\mu = \frac{m}{L} = \frac{3.9\times 10^{-3}}{0.65} = \SI{6.0e-3}{kg.m^{-1}}.
\]
\textbf{Step 2: wave speed.}
\[
v = \sqrt{\frac{T}{\mu}} = \sqrt{\frac{96}{6.0\times 10^{-3}}}
= \sqrt{1.6\times 10^{4}} \approx \SI{1.3e2}{m.s^{-1}}.
\]
The pulse crosses the whole string in about $L/v = 0.65/126 \approx \SI{5.2}{ms}$, which is why the eye cannot follow it.
\end{exoresolu}

\begin{exoresolu}[Finding the tension of a piano string]
A steel piano string has linear density $\mu = \SI{5.0e-3}{kg.m^{-1}}$. A transverse wave is measured to travel along it at $v = \SI{2.0e2}{m.s^{-1}}$. Find the tension in the string.
\tcblower
Square the wave-speed formula and solve for $T$:
\[
v = \sqrt{\frac{T}{\mu}} \;\Rightarrow\; v^{2} = \frac{T}{\mu} \;\Rightarrow\; T = \mu v^{2}.
\]
Numerically:
\[
T = 5.0\times 10^{-3}\times(2.0\times 10^{2})^{2}
= 5.0\times 10^{-3}\times 4.0\times 10^{4}
= \SI{2.0e2}{N}.
\]
The string is pulled with a force of $\SI{200}{N}$, roughly the weight of a $\SI{20}{kg}$ mass. A full piano frame must withstand the sum of such tensions from more than two hundred strings, which is why it is made of cast iron.
\end{exoresolu}

\begin{exoresolu}[Finding the linear density of a wire]
A wire under a tension of $\SI{50}{N}$ carries transverse waves at $\SI{25}{m.s^{-1}}$. Calculate its linear density, and then the mass of a $\SI{2.0}{m}$ length of this wire.
\tcblower
From $v = \sqrt{T/\mu}$ we obtain $\mu = T/v^{2}$:
\[
\mu = \frac{50}{25^{2}} = \frac{50}{625} = \SI{8.0e-2}{kg.m^{-1}}.
\]
Mass of a $\SI{2.0}{m}$ length:
\[
m = \mu L = 8.0\times 10^{-2}\times 2.0 = \SI{0.16}{kg}.
\]
\end{exoresolu}

\begin{exoresolu}[Frequency of a vibrating string]
The guitar string of the first worked example ($L = \SI{0.65}{m}$, $\mu = \SI{6.0e-3}{kg.m^{-1}}$, $T = \SI{96}{N}$) vibrates in its fundamental mode, for which $\lambda = 2L$. Calculate its fundamental frequency.
\tcblower
We already found $v = \SI{1.3e2}{m.s^{-1}}$ (more precisely $\SI{126.5}{m.s^{-1}}$). The wavelength is
\[
\lambda = 2L = 2\times 0.65 = \SI{1.30}{m}.
\]
Hence
\[
f = \frac{v}{\lambda} = \frac{126.5}{1.30} \approx \SI{97}{Hz}.
\]
This is close to the note G$_2$ ($\SI{98}{Hz}$), the lowest string of a standard guitar: the physics matches the music.
\end{exoresolu}

\begin{exoresolu}[Proportional reasoning: retuning a string]
A string sounds a fundamental frequency of $\SI{110}{Hz}$ under a tension of $\SI{80}{N}$. The musician wants the same string to sound $\SI{165}{Hz}$ without changing its length. What tension is required?
\tcblower
At fixed $L$ and $\mu$, the fundamental frequency obeys $f_1 \propto \sqrt{T}$. Therefore
\[
\frac{T'}{T} = \left(\frac{f'_1}{f_1}\right)^{2}
= \left(\frac{165}{110}\right)^{2}
= (1.5)^{2} = 2.25.
\]
Hence
\[
T' = 2.25\times 80 = \SI{1.8e2}{N}.
\]
Raising the pitch by a factor $1.5$ requires \emph{more than doubling} the tension, because the frequency depends only on the \emph{square root} of $T$. This is why tuning pegs must be turned with care.
\end{exoresolu}

\courssub{Common Mistakes and Key Points}

\begin{attention}[Using the mass $m$ instead of the linear density $\mu$]
The most frequent error at the GCE is to write $v = \sqrt{T/m}$, inserting the \emph{total mass} of the string. This is wrong, and dimensional analysis exposes it immediately:
\[
\left[\frac{T}{m}\right] = \frac{\mathrm{M\,L\,T^{-2}}}{\mathrm{M}} = \mathrm{L\,T^{-2}} \neq \mathrm{L^{2}\,T^{-2}}.
\]
The formula needs the mass \emph{per unit length}, $\mu = m/L$. Other traps to avoid:
\begin{itemize}
\item forgetting to convert $\mu$ given in $\mathrm{g\,m^{-1}}$ into $\mathrm{kg\,m^{-1}}$ (divide by $1000$);
\item forgetting the square root, or squaring only one side when solving for $T$ (remember $T = \mu v^{2}$);
\item believing that increasing the frequency increases the speed: for a fixed string under fixed tension, $v$ is unchanged; only $\lambda$ decreases;
\item mixing up the length $L$ of the string with the wavelength $\lambda$: in the fundamental mode $\lambda = 2L$, not $L$;
\item in proportional-reasoning questions, forgetting that $v$ and $f_1$ vary as $\sqrt{T}$, so a factor of $n$ in frequency requires a factor of $n^{2}$ in tension.
\end{itemize}
\end{attention}

\begin{exercice}[Checking a laboratory measurement]
In a school experiment, a rubber cord of length $\SI{1.5}{m}$ and mass $\SI{0.030}{kg}$ is stretched by a force of $\SI{20}{N}$. (a) Calculate the speed of transverse waves on the cord. (b) The cord is now stretched to $\SI{80}{N}$ without its linear density changing appreciably. By what factor is the speed multiplied? (c) State the fundamental frequency of the cord in case (b).
\end{exercice}

\begin{corrige}[Exercise 1]
(a) $\mu = m/L = 0.030/1.5 = \SI{2.0e-2}{kg.m^{-1}}$, so $v = \sqrt{20/(2.0\times 10^{-2})} = \sqrt{1000} \approx \SI{32}{m.s^{-1}}$.
(b) Since $v \propto \sqrt{T}$ at fixed $\mu$, quadrupling $T$ from $\SI{20}{N}$ to $\SI{80}{N}$ doubles the speed: $v' = 2v \approx \SI{63}{m.s^{-1}}$.
(c) Fundamental mode: $\lambda = 2L = \SI{3.0}{m}$, so $f = v'/\lambda = 63.2/3.0 \approx \SI{21}{Hz}$.
\end{corrige}

\begin{exercice}[Two strings on the same instrument]
Two strings of the same length $\SI{0.60}{m}$ are fixed on an instrument. String A has $\mu_A = \SI{4.0e-3}{kg.m^{-1}}$ and is under tension $T_A = \SI{90}{N}$. String B must sound a fundamental note one octave \emph{below} string A (i.e. at half the frequency) under the same tension. (a) Calculate the fundamental frequency of string A. (b) Find the linear density $\mu_B$ required for string B.
\end{exercice}

\begin{corrige}[Exercise 2]
(a) $v_A = \sqrt{90/(4.0\times 10^{-3})} = \sqrt{2.25\times 10^{4}} = \SI{1.5e2}{m.s^{-1}}$; with $\lambda = 2L = \SI{1.20}{m}$, $f_A = 150/1.20 = \SI{125}{Hz}$.
(b) At fixed $L$ and $T$, $f_1 \propto 1/\sqrt{\mu}$. Halving the frequency therefore requires multiplying $\mu$ by $2^{2} = 4$: $\mu_B = 4\times 4.0\times 10^{-3} = \SI{1.6e-2}{kg.m^{-1}}$. (Check: $v_B = \sqrt{90/(1.6\times 10^{-2})} = \SI{75}{m.s^{-1}}$, so $f_B = 75/1.20 = \SI{62.5}{Hz} = f_A/2$.)
\end{corrige}

\begin{aretenir}[Key points]
\begin{itemize}
\item The linear density of a string is $\mu = m/L$, in $\mathrm{kg\,m^{-1}}$.
\item The speed of a transverse wave on a taut string is $v = \sqrt{T/\mu}$, with $T$ in newtons and $v$ in $\mathrm{m\,s^{-1}}$.
\item Dimensional analysis gives the form of the law: $[T/\mu] = \mathrm{L^{2}\,T^{-2}}$, whose square root is a speed.
\item The speed depends only on the medium ($T$ and $\mu$), never on the frequency or the amplitude.
\item Combined with $v = f\lambda$, the formula explains how stringed instruments are tuned: $f_1 = \dfrac{1}{2L}\sqrt{\dfrac{T}{\mu}}$.
\item At fixed $\mu$, $v \propto \sqrt{T}$; at fixed $L$ and $\mu$, $f_1 \propto \sqrt{T}$; at fixed $L$ and $T$, $f_1 \propto 1/\sqrt{\mu}$.
\item Always check dimensions and convert all quantities to SI units before calculating.
\end{itemize}
\end{aretenir}

\begin{savaistu}[Why are bass piano strings wound with copper?]
To produce a low note at a reasonable tension and length, a string needs a \emph{large} linear density $\mu$, since $f_1 = \frac{1}{2L}\sqrt{T/\mu}$. A plain steel wire thick enough would be too stiff to vibrate musically. Piano makers therefore wrap a thin steel core with copper wire: this raises $\mu$ (and hence lowers the pitch) while keeping the string flexible. The same trick is used for the bass strings of guitars, including those played by Cameroonian makossa and bikutsi guitarists.
\end{savaistu}

\coursec{Experimental Study and Applications: The Sonometer and Vibrating Strings}

In the previous section we established, both by dimensional analysis and by a proper derivation, that the speed of a transverse wave on a taut string is $v = \sqrt{T/\mu}$, where $T$ is the tension and $\mu$ the mass per unit length. A formula in physics is only trustworthy when experiment confirms it. In this section we meet the classic instrument designed for exactly that purpose --- the \cle{sonometer} --- and we use it to verify the law, to study the \cle{stationary waves} that a stretched string can sustain, and to understand how musicians in Yaound\'e or anywhere else tune a guitar, a piano or an \emph{mvet}.

\courssub{The Sonometer: Description and Use}

\begin{definition}[The sonometer]
A \cle{sonometer} is a hollow wooden sounding box on which a thin wire is stretched between two \cle{bridges}. One end of the wire is fixed; the other passes over a pulley and carries a \cle{hanger} on which slotted masses are placed, so that the tension $T$ in the wire equals the total weight supported, $T = Mg$. A \cle{movable bridge} sliding under the wire allows the vibrating length $L$ to be adjusted and measured on a graduated scale fixed to the box.
\end{definition}

The hollow box is not a luxury: it acts as a resonator, amplifying the faint sound of the vibrating wire so that the experimenter can hear it clearly. The wire is set into vibration either by plucking it at its centre or by pressing against it a small \emph{rider} (a tiny paper or wire loop) which is thrown off when the vibration amplitude becomes large --- a convenient way of detecting resonance.

\begin{popfigure}
\begin{tikzpicture}[scale=1.0, every node/.style={font=\small}]
% sounding box
\draw[fill=popblueL, draw=popink, thick] (0,0) rectangle (10,1.2);
\draw[popmuted] (0.3,0.6) -- (9.7,0.6);
\node[popmuted] at (5,0.3) {hollow sounding box};
% scale
\draw[popdark] (0.4,1.05) -- (9.2,1.05);
\foreach \x in {0.4,1.4,2.4,3.4,4.4,5.4,6.4,7.4,8.4,9.2} \draw[popdark] (\x,1.0) -- (\x,1.1);
% fixed bridge
\draw[fill=poporange, draw=popdark] (0.7,1.2) -- (1.1,1.2) -- (0.9,1.7) -- cycle;
\node[above, popdark] at (0.9,1.75) {fixed bridge};
% movable bridge
\draw[fill=popgreen, draw=popdark] (6.3,1.2) -- (6.7,1.2) -- (6.5,1.7) -- cycle;
\node[above, popdark] at (6.5,1.75) {movable bridge};
% wire
\draw[very thick, popink] (-0.6,1.7) -- (10.2,1.7);
% fixed end
\draw[fill=popdark] (-0.6,1.7) circle (0.07);
\node[left, popdark] at (-0.6,1.7) {fixed end};
% pulley
\draw[fill=white, draw=popdark, thick] (10.6,1.7) circle (0.4);
\draw[fill=popdark] (10.6,1.7) circle (0.05);
% string over pulley and hanger
\draw[very thick, popink] (10.2,1.7) -- (10.6,2.1);
\draw[very thick, popink] (10.6,2.1) arc (90:-90:0.4);
\draw[very thick, popink] (10.6,1.3) -- (10.6,0.4);
\draw[fill=popgold, draw=popdark] (10.25,0.1) rectangle (10.95,0.4);
\node[right, popdark] at (11.0,0.25) {hanger + masses};
% vibrating length
\draw[<->, poppurple, thick] (0.9,2.35) -- (6.5,2.35);
\node[poppurple, above] at (3.7,2.35) {vibrating length $L$};
% tension label
\node[popdark] at (8.6,2.0) {wire under tension $T = Mg$};
\end{tikzpicture}
\legende{The sonometer: a wire stretched over two bridges on a hollow box, tensioned by hanging masses; the movable bridge sets the vibrating length $L$.}
\end{popfigure}

\courssub{Stationary Waves on a String Fixed at Both Ends}

When a stretched string is plucked, waves travel along it in both directions, reflect at the fixed ends, and superpose. Because both ends are fixed, they must be \cle{nodes} (points of zero displacement). Only certain wavelengths can ``fit'' between the two ends: these are the \cle{modes of vibration} of the string.

\begin{propriete}[Stationary-wave condition for a string fixed at both ends]
A string of length $L$ fixed at both ends can sustain a stationary wave only if an integer number of half-wavelengths fits into its length:
\[ L = n\,\frac{\lambda_n}{2}, \qquad n = 1, 2, 3, \dots \]
so the allowed wavelengths are $\lambda_n = \dfrac{2L}{n}$ and the allowed frequencies are
\[ f_n = \frac{v}{\lambda_n} = \frac{nv}{2L} = n f_1 . \]
The mode $n = 1$ is the \cle{fundamental} (first harmonic); the modes $n = 2, 3, \dots$ are the higher \cle{harmonics} (overtones). The derivation (substituting $\lambda_n = 2L/n$ into $f_n = v/\lambda_n$) is short and examinable.
\end{propriete}

\begin{definition}[Fundamental frequency and harmonics]
The \cle{fundamental frequency} $f_1$ of a stretched string is the lowest frequency at which it can vibrate; the string then vibrates in a single loop with one antinode at the centre. The \cle{harmonics} are the frequencies $f_n = n f_1$; the $n$-th harmonic has $n$ loops, $n$ antinodes and $n+1$ nodes (including the two fixed ends).
\end{definition}

\begin{popfigure}
\begin{tikzpicture}[scale=1.0, every node/.style={font=\small}]
% Fundamental
\begin{scope}
\draw[popdark, thick] (0,0) -- (6,0);
\draw[fill=popdark] (0,0) circle (0.06); \draw[fill=popdark] (6,0) circle (0.06);
\draw[very thick, popblue, domain=0:6, samples=60] plot (\x, {0.8*sin(30*\x)});
\draw[dashed, popmuted, domain=0:6, samples=60] plot (\x, {-0.8*sin(30*\x)});
\node[popink] at (3,-1.5) {Fundamental: $n=1$, $\lambda_1 = 2L$};
\node[popink] at (3,-2.0) {$f_1 = \dfrac{v}{2L}$};
\node[popdark, above] at (3,0.85) {1 loop};
\end{scope}
% 2nd harmonic
\begin{scope}[yshift=-3.2cm]
\draw[popdark, thick] (0,0) -- (6,0);
\draw[fill=popdark] (0,0) circle (0.06); \draw[fill=popdark] (6,0) circle (0.06);
\draw[very thick, poporange, domain=0:6, samples=80] plot (\x, {0.8*sin(60*\x)});
\draw[dashed, popmuted, domain=0:6, samples=80] plot (\x, {-0.8*sin(60*\x)});
\draw[fill=poppurple] (3,0) circle (0.06);
\node[poppurple, above] at (3,0.1) {N};
\node[popink] at (3,-1.5) {2nd harmonic: $n=2$, $\lambda_2 = L$};
\node[popink] at (3,-2.0) {$f_2 = 2f_1$};
\node[popdark, above] at (3,0.85) {2 loops};
\end{scope}
% 3rd harmonic
\begin{scope}[yshift=-6.4cm]
\draw[popdark, thick] (0,0) -- (6,0);
\draw[fill=popdark] (0,0) circle (0.06); \draw[fill=popdark] (6,0) circle (0.06);
\draw[very thick, popgreen, domain=0:6, samples=100] plot (\x, {0.8*sin(90*\x)});
\draw[dashed, popmuted, domain=0:6, samples=100] plot (\x, {-0.8*sin(90*\x)});
\draw[fill=poppurple] (2,0) circle (0.06); \draw[fill=poppurple] (4,0) circle (0.06);
\node[popink] at (3,-1.5) {3rd harmonic: $n=3$, $\lambda_3 = \dfrac{2L}{3}$};
\node[popink] at (3,-2.0) {$f_3 = 3f_1$};
\node[popdark, above] at (3,0.85) {3 loops};
\end{scope}
\end{tikzpicture}
\legende{The first three stationary-wave modes of a string fixed at both ends. N marks an interior node; the dashed curve shows the string half a period later.}
\end{popfigure}

Combining the stationary-wave condition with the wave-speed formula $v = \sqrt{T/\mu}$ gives the central result of this section.

\begin{propriete}[Fundamental frequency of a stretched string]
For a string of length $L$ fixed at both ends, under tension $T$, of mass per unit length $\mu$, the fundamental frequency is
\[ f_1 = \frac{1}{2L}\sqrt{\frac{T}{\mu}} . \]
This shows directly how a player controls the pitch: $f_1$ increases when $L$ decreases, when $T$ increases, or when $\mu$ decreases. (The derivation --- substituting $\lambda_1 = 2L$ into $f_1 = v/\lambda_1$ with $v = \sqrt{T/\mu}$ --- is short and examinable.)
\end{propriete}

\begin{attention}[The fundamental wavelength is $\lambda = 2L$, not $L$]
The most frequent error in GCE scripts is to write $\lambda_1 = L$ for the fundamental mode. Remember: the fundamental has \textbf{one loop}, and one loop is \textbf{half} a wavelength (the distance between two consecutive nodes is $\lambda/2$). Hence $\lambda_1 = 2L$ and $f_1 = v/(2L)$. Writing $\lambda_1 = L$ doubles your frequency and destroys the rest of the calculation.
\end{attention}

\courssub{Experiment 1: Verifying $v \propto \sqrt{T}$ at Constant $\mu$}

\begin{methode}[Protocol: verifying $v = \sqrt{T/\mu}$ with a sonometer]
\begin{enumerate}
\item Set up the sonometer with a single wire of measured mass per unit length $\mu$ (obtained by weighing a known length of the same wire: $\mu = m/\ell$).
\item Place a mass $M$ on the hanger so that $T = Mg$. Keep the same wire throughout, so $\mu$ is constant.
\item Fix the vibrating length $L$ with the movable bridge. Pluck the wire gently at its centre and compare its note with a reference (tuning fork or frequency generator); adjust $L$ until the wire resonates at the chosen reference frequency $f$ (a paper rider at the centre is thrown off at resonance).
\item Record $T$ and the resonant length $L$. Since the wire vibrates in its fundamental mode, $v = f\lambda_1 = 2fL$.
\item Repeat for at least five different masses $M$, keeping $f$ fixed.
\item Plot $L$ against $\sqrt{T}$ (equivalently $v$ against $\sqrt{T}$). A straight line through the origin confirms $v \propto \sqrt{T}$. The gradient of the $v$ versus $\sqrt{T}$ graph equals $1/\sqrt{\mu}$, from which $\mu$ can be deduced and compared with the directly measured value.
\end{enumerate}
\end{methode}

\begin{exemplebox}[Numerical illustration]
With $f = \SI{256}{Hz}$ fixed, a student finds resonance at $L = \SI{0.500}{m}$ for $M = \SI{4.0}{kg}$. Then $T = Mg = 4.0 \times 9.81 = \SI{39.2}{N}$ and $v = 2fL = 2 \times 256 \times 0.500 = \SI{256}{m.s^{-1}}$. Doubling the mass to $M = \SI{8.0}{kg}$ should multiply $v$ by $\sqrt{2}$, so the predicted resonant length is $L' = 0.500\sqrt{2} = \SI{0.707}{m}$ --- exactly what the experiment shows within a few per cent.
\end{exemplebox}

\courssub{Experiment 2: Verifying $v \propto 1/\sqrt{\mu}$ at Constant $T$}

Now keep the hanging mass fixed (constant $T$) and use, in turn, several wires of the same material but different diameters, or of different materials (steel, brass, copper). For each wire, measure the resonant length $L$ at the same reference frequency $f$, so $v = 2fL$.

Since $\mu = \rho \pi d^{2}/4$ for a wire of density $\rho$ and diameter $d$, the theory predicts
\[ v = \sqrt{\frac{T}{\mu}} \quad\Longrightarrow\quad v \propto \frac{1}{\sqrt{\mu}} \propto \frac{1}{d} \text{ (same material).} \]
Plotting $v$ against $1/\sqrt{\mu}$ (or $L$ against $1/d$ for wires of one material) gives a straight line through the origin, confirming the law. A thicker wire vibrates more slowly --- this is why the bass strings of a guitar are thick and wound.

\begin{propriete}[Effect of tension on speed and frequency]
Because $v = \sqrt{T/\mu}$ and $f_1 = v/(2L)$, both the wave speed and the fundamental frequency vary as $\sqrt{T}$. \textbf{Doubling the tension multiplies $v$ and $f_1$ by $\sqrt{2} \approx 1.41$, not by 2.} To double the frequency of a string (raise it by one octave) at fixed $L$ and $\mu$, the tension must be multiplied by \textbf{four}.
\end{propriete}

\courssub{Applications: Tuning String Instruments}

Every string instrument is an applied sonometer. The formula $f_1 = \dfrac{1}{2L}\sqrt{\dfrac{T}{\mu}}$ explains the three gestures of any musician:

\begin{itemize}
\item \textbf{Tightening a peg} increases $T$, so $f_1$ rises: the note becomes sharper. Loosening lowers the pitch. A guitarist tunes each string this way.
\item \textbf{Pressing a finger on the fingerboard} shortens the vibrating length $L$, raising the frequency: this is how different notes are produced on one string of a guitar or of the \emph{mvet}, the traditional harp-zither of the Beti people.
\item \textbf{Choosing the string}: bass notes require low $f_1$, hence large $\mu$ --- thick or wound strings. Treble strings are thin. In a piano, the bass strings are long, thick and copper-wound, while the treble strings are short and thin.
\end{itemize}

\begin{savaistu}[The mvet and the physics of strings]
The \emph{mvet}, emblematic instrument of the Fang-Beti epic tradition of the South and Centre regions of Cameroon, carries several strings stretched over a resonating body, exactly like a multi-wire sonometer. The player shortens the vibrating length with the fingers and selects strings of different thicknesses to cover a wide range of pitches --- a living illustration of $f_1 = \frac{1}{2L}\sqrt{T/\mu}$, centuries before the formula was written down.
\end{savaistu}

\courssub{Data Treatment and GCE-Style Problem Solving}

\begin{methode}[Solving problems combining $v = f\lambda$, $v = \sqrt{T/\mu}$ and $L = n\lambda/2$]
\begin{enumerate}
\item Identify the mode: fundamental ($\lambda = 2L$) or $n$-th harmonic ($\lambda = 2L/n$). \textbf{Draw the loop pattern if in doubt.}
\item Write $v = f\lambda$ using the wavelength of the correct mode.
\item Write $v = \sqrt{T/\mu}$ with $T = Mg$ if masses hang from the wire.
\item Equate or combine the two expressions; solve for the unknown, keeping units in SI (lengths in metres, $\mu$ in $\mathrm{kg\,m^{-1}}$, $T$ in newtons).
\item For graph questions: rearrange the formula so that the plotted quantities give a straight line, then read the gradient and interpret it physically (e.g. gradient of $f_1$ versus $1/L$ equals $\frac{1}{2}\sqrt{T/\mu}$).
\end{enumerate}
\end{methode}

\begin{exoresolu}[Fundamental frequency of a sonometer wire]
A sonometer wire of length $L = \SI{0.80}{m}$ between the bridges has mass per unit length $\mu = \SI{1.2e-3}{kg.m^{-1}}$. It is tensioned by a mass $M = \SI{5.0}{kg}$ (take $g = \SI{9.81}{m.s^{-2}}$).
\begin{enumerate}
\item Calculate the speed of transverse waves on the wire.
\item Calculate the fundamental frequency and the frequency of the second harmonic.
\item The tension is doubled. State the new fundamental frequency.
\end{enumerate}
\tcblower
\textbf{Solution.}
\begin{enumerate}
\item Tension: $T = Mg = 5.0 \times 9.81 = \SI{49.05}{N}$.
\[ v = \sqrt{\frac{T}{\mu}} = \sqrt{\frac{49.05}{1.2\times10^{-3}}} = \sqrt{4.09\times10^{4}} \approx \SI{202}{m.s^{-1}}. \]
\item Fundamental mode: $\lambda_1 = 2L = \SI{1.60}{m}$.
\[ f_1 = \frac{v}{2L} = \frac{202}{1.60} \approx \SI{126}{Hz}. \]
Second harmonic: $f_2 = 2f_1 \approx \SI{253}{Hz}$ (the string vibrates in two loops, $\lambda_2 = L = \SI{0.80}{m}$).
\item Doubling $T$ multiplies $f_1$ by $\sqrt{2}$: $f_1' = 126 \times 1.414 \approx \SI{179}{Hz}$ --- \emph{not} $\SI{252}{Hz}$.
\end{enumerate}
\end{exoresolu}

\begin{exoresolu}[Deducing $\mu$ from a graph]
In a sonometer experiment at fixed frequency $f = \SI{200}{Hz}$, the resonant length $L$ is measured for various tensions $T$. The graph of $L$ (vertical axis) against $\sqrt{T}$ is a straight line through the origin of gradient $k = \SI{0.045}{m.N^{-1/2}}$. Deduce the mass per unit length $\mu$ of the wire.
\tcblower
\textbf{Solution.} At resonance in the fundamental mode, $v = 2fL$ and $v = \sqrt{T/\mu}$, so
\[ 2fL = \sqrt{\frac{T}{\mu}} \quad\Longrightarrow\quad L = \frac{1}{2f\sqrt{\mu}}\,\sqrt{T}. \]
The graph of $L$ versus $\sqrt{T}$ has gradient $k = \dfrac{1}{2f\sqrt{\mu}}$. Hence
\[ \sqrt{\mu} = \frac{1}{2fk} = \frac{1}{2 \times 200 \times 0.045} = \frac{1}{18} = \SI{0.0556}{kg^{1/2}.m^{-1/2}}, \]
\[ \mu = (0.0556)^2 \approx \SI{3.1e-3}{kg.m^{-1}}. \]
\end{exoresolu}

\begin{exercice}[Vibrating strings]
A steel guitar string has length $\SI{0.65}{m}$ and mass per unit length $\SI{4.0e-3}{kg.m^{-1}}$.
\begin{enumerate}
\item What tension is required for a fundamental frequency of $\SI{330}{Hz}$?
\item With this tension, what are the wavelengths and frequencies of the second and third harmonics?
\item The player presses the string so that only $\SI{0.50}{m}$ vibrates. What is the new fundamental frequency?
\end{enumerate}
\end{exercice}

\begin{corrige}[Exercise 1]
\begin{enumerate}
\item $f_1 = \dfrac{1}{2L}\sqrt{T/\mu} \Rightarrow T = 4L^{2}f_1^{2}\mu = 4(0.65)^2(330)^2(4.0\times10^{-3}) \approx \SI{736}{N}$.
\item $f_2 = 2f_1 = \SI{660}{Hz}$, $\lambda_2 = L = \SI{0.65}{m}$; $f_3 = 3f_1 = \SI{990}{Hz}$, $\lambda_3 = 2L/3 = \SI{0.43}{m}$.
\item $f_1' = f_1 \times \dfrac{L}{L'} = 330 \times \dfrac{0.65}{0.50} = \SI{429}{Hz}$ (frequency is inversely proportional to length at fixed $T$ and $\mu$).
\end{enumerate}
\end{corrige}

\begin{aretenir}[Key points]
\begin{itemize}
\item The \cle{sonometer} verifies $v = \sqrt{T/\mu}$: vary $T$ at fixed $\mu$ (straight line of $v$ versus $\sqrt{T}$), then vary $\mu$ at fixed $T$ (straight line of $v$ versus $1/\sqrt{\mu}$).
\item A string fixed at both ends vibrates only in modes satisfying $L = n\lambda/2$; the fundamental has $\lambda_1 = 2L$.
\item Fundamental frequency: $f_1 = \dfrac{1}{2L}\sqrt{\dfrac{T}{\mu}}$; harmonics: $f_n = nf_1$.
\item Doubling $T$ multiplies $v$ and $f_1$ by $\sqrt{2}$; doubling the frequency requires four times the tension.
\item Thick strings $\Rightarrow$ large $\mu$ $\Rightarrow$ low notes; tightening a peg $\Rightarrow$ larger $T$ $\Rightarrow$ sharper note; shortening $L$ $\Rightarrow$ higher note.
\end{itemize}
\end{aretenir}

\newpage

\coursec{Integration activity}

\courssub{Aims of the activity}

By the end of this integration activity, you should be able to:

\begin{itemize}
\item Apply the relation $v=\sqrt{T/\mu}$ to determine the speed of a transverse wave on a taut string from its tension $T$ and linear density $\mu$, and conversely to find the tension needed for a given speed.
\item Combine $v=f\lambda$ with the stationary-wave condition for a string fixed at both ends, $\lambda_1=2L$ for the fundamental mode, to design a string that produces a required musical note.
\item Predict quantitatively how the fundamental frequency changes when the tension or the linear density of the string is modified, using the proportionality $f_1\propto\sqrt{T}/\sqrt{\mu}$.
\item Use the rope-and-slit model to explain the \cle{polarization} of a transverse wave, and interpret the darkening observed when two polaroid sheets are rotated relative to each other.
\item Present a reasoned scientific argument, with correct SI units and sensible significant figures, to the class.
\end{itemize}

This activity brings together the two lessons of the chapter: the \cle{polarization} of transverse waves and the \cle{speed of transverse waves} on taut strings and wires. It shows that both topics rest on one central idea: the \emph{transverse nature} of the waves involved. On a string, the transverse displacement propagates at a speed fixed by the mechanical properties of the medium ($T$ and $\mu$); in polarization, it is precisely the transverse direction of vibration that is selected by a slit or a polaroid sheet.

\courssub{Situation}

\textbf{Tuning the Mvet: Designing a String for a Given Note.}\par

During the cultural week of your school in Yaoundé, the music club decides to build a simplified version of the \emph{mvet}, the traditional harp-zither of the forest regions of Cameroon. In the traditional instrument, the player stretches fibre strings along a wooden stick fitted with resonating calabashes, and tunes each string by adjusting its tension until the desired note is obtained.

The club asks your physics group to do the scientific design work. The specifications are the following:

\begin{itemize}
\item The vibrating part of the string is fixed at both ends and has length $L=\SI{65}{cm}=\SI{0.65}{m}$.
\item The string is made of plant fibre of linear density (mass per unit length) $\mu=\SI{1.2}{g.m^{-1}}=\SI{1.2e-3}{kg.m^{-1}}$.
\item When plucked at its centre, the string must emit as its \emph{fundamental} note the standard pitch $f_1=\SI{220}{Hz}$ (the note A used to tune orchestras).
\end{itemize}

The club also wants to understand two practical questions before buying materials:

\begin{itemize}
\item What happens to the note if, after some days of playing, the musician tightens the string so that its tension \emph{increases by} $20\,\%$?
\item What happens if, instead of the original fibre, a fibre string \emph{twice as thick} is used (same material, so twice the linear density), stretched under the \emph{same} tension?
\end{itemize}

Finally, during the open day, your group must run a small optics stand. You are given two square \emph{polaroid sheets} taken from old sunglasses. Visitors notice that when the two sheets are held one behind the other towards the sunlight and one sheet is slowly rotated, the transmitted light becomes progressively dimmer, until the view is almost completely dark when the sheets are ``crossed'' at $90^{\circ}$. You must explain this observation using the rope-and-slit model studied in class, and state clearly what it proves about the nature of light.

\begin{attention}[Units before anything else]
Convert all data to SI units \emph{before} any calculation: $L$ in metres, $\mu$ in $\mathrm{kg\,m^{-1}}$, $T$ in newtons, $f$ in hertz. Most errors in this type of problem come from forgetting to convert $\mathrm{g\,m^{-1}}$ into $\mathrm{kg\,m^{-1}}$, or $\mathrm{cm}$ into $\mathrm{m}$. A second frequent error is to confuse the wavelength $\lambda_1=2L$ with the length $L$ itself.
\end{attention}

\begin{popfigure}
\begin{center}
\begin{tikzpicture}[scale=1.0]
% Rope and two slits
\draw[popink,thick,domain=0:3.2,samples=60] plot (\x,{0.45*sin(360*\x/1.6)});
\draw[popink,thick] (3.2,0) -- (4.0,0);
% Slit 1 (vertical)
\fill[popdark] (4.0,-1.0) rectangle (4.25,-0.28);
\fill[popdark] (4.0,0.28) rectangle (4.25,1.0);
\node[popdark,font=\small] at (4.12,1.25) {slit 1};
% Rope between slits
\draw[popink,thick,domain=4.25:6.8,samples=60] plot (\x,{0.45*sin(360*(\x-1.05)/1.6)});
\draw[popink,thick] (6.8,0) -- (7.6,0);
% Slit 2 (horizontal)
\fill[poporange] (7.6,-1.0) rectangle (8.6,-0.72);
\fill[poporange] (7.6,0.72) rectangle (8.6,1.0);
\node[poporange,font=\small] at (8.1,1.25) {slit 2 (rotated $90^{\circ}$)};
% Rope after slit 2: no vibration
\draw[popink,thick] (8.6,0) -- (10.6,0);
\node[popmuted,font=\small,align=center,text width=3.2cm] at (9.7,-0.6) {no vibration beyond slit 2};
\node[popmuted,font=\small,align=center,text width=3.4cm] at (1.6,-1.1) {plane polarized transverse wave};
\end{tikzpicture}
\end{center}
\legende{Rope-and-slit model: the first vertical slit selects a vertical plane of vibration (plane polarization); the second slit, rotated by $90^{\circ}$, blocks the vibration completely.}
\end{popfigure}

\courssub{Tasks}

Work in groups of three or four. Answer the questions in order; each answer will be used in the next question.

\begin{enumerate}
\item \textbf{Stationary-wave condition.} The string is fixed at both ends. Sketch the fundamental mode of vibration and justify that the wavelength of the fundamental is $\lambda_1=2L$. Calculate $\lambda_1$ numerically.
\item \textbf{Required wave speed.} Using $v=f\lambda$, calculate the speed that transverse waves must have on the string for the fundamental note to be $f_1=\SI{220}{Hz}$.
\item \textbf{Required tension.} Using $v=\sqrt{T/\mu}$, show that the tension to be applied is $T=\mu v^{2}$, and calculate its value in newtons. Compare this tension with the weight of a stone of mass $m$: what mass hung from the string would produce the same tension (take $g=\SI{9.81}{m.s^{-2}}$)?
\item \textbf{Effect of a $20\,\%$ increase of tension.} The tension becomes $T'=1.20\,T$, with $L$ and $\mu$ unchanged. Show that the new fundamental frequency is $f_1'=f_1\sqrt{1.20}$ and calculate it. Does the note sound higher or lower? Is the change audible (a change of more than about $3\,\%$ is easily heard)?
\item \textbf{Effect of doubling the thickness.} A string of the same material but twice as thick has $\mu''=2\mu$ and is stretched under the same tension $T$. Show that $f_1''=f_1/\sqrt{2}$ and calculate it. Why do instrument makers use \emph{thicker} strings for the \emph{low} notes?
\item \textbf{Polarization demonstration.} (a) Describe the rope-and-slit experiment: a rope passes through two parallel slits; what happens to the rope's transverse vibration when the second slit is rotated by $90^{\circ}$? (b) By analogy, explain why the light transmitted by two polaroid sheets becomes very weak when their transmission axes are crossed at $90^{\circ}$. (c) State precisely what this experiment proves about light, and why the same experiment performed with sound waves in air would show no such effect.
\item \textbf{Synthesis.} Prepare a five-minute presentation: your calculations for questions 1 to 5 (with units and significant figures), and a live demonstration for question 6, ending with one sentence that links the two parts of the activity to the transverse nature of waves.
\end{enumerate}

\courssub{Model answer}

\textbf{1. Stationary-wave condition.}\par
The string is fixed at both ends, so each end must be a \cle{node} of vibration (a point that never moves). The simplest stationary pattern fitting between two nodes is a single loop, with one \cle{antinode} at the centre. The length of the string then equals half a wavelength, since one full wavelength contains two loops:
\[ L=\frac{\lambda_1}{2}\quad\Longrightarrow\quad \lambda_1=2L=2\times 0.65=\SI{1.30}{m}. \]

\textbf{2. Required wave speed.}\par
From the general wave relation $v=f\lambda$ applied to the fundamental mode:
\[ v=f_1\lambda_1=220\times 1.30=\SI{286}{m.s^{-1}}. \]
Transverse waves must therefore travel along the fibre at about $\SI{286}{m.s^{-1}}$.

\textbf{3. Required tension.}\par
Starting from $v=\sqrt{T/\mu}$ and squaring both sides:
\[ v^{2}=\frac{T}{\mu}\quad\Longrightarrow\quad T=\mu v^{2}=1.2\times 10^{-3}\times(286)^{2}. \]
\[ T=1.2\times 10^{-3}\times 8.18\times 10^{4}\approx \SI{98}{N}. \]
A mass $m$ hung vertically from the string would produce a tension $T=mg$, so
\[ m=\frac{T}{g}=\frac{98}{9.81}\approx \SI{10}{kg}. \]
The string must therefore be tightened with a force equivalent to the weight of a $10\ \mathrm{kg}$ bag — which is exactly what the tuning pegs of the instrument are designed to do.

\textbf{4. Effect of a $20\,\%$ increase of tension.}\par
Since $f_1=\dfrac{v}{2L}=\dfrac{1}{2L}\sqrt{\dfrac{T}{\mu}}$, with $L$ and $\mu$ fixed the frequency is proportional to $\sqrt{T}$. Hence
\[ f_1'=f_1\sqrt{\frac{T'}{T}}=220\times\sqrt{1.20}=220\times 1.095\approx \SI{241}{Hz}. \]
The note rises by about $9.5\,\%$: it sounds clearly \emph{higher} (roughly a semitone sharper). This is why a mvet must be retuned whenever the strings are tightened or when humidity changes the tension of the fibres.

\textbf{5. Effect of doubling the thickness.}\par
With $T$ and $L$ unchanged, $f_1\propto 1/\sqrt{\mu}$:
\[ f_1''=f_1\sqrt{\frac{\mu}{\mu''}}=\frac{220}{\sqrt{2}}\approx \SI{156}{Hz}. \]
The note becomes \emph{lower}. Thicker strings have a larger linear density, so waves travel more slowly on them ($v=\sqrt{T/\mu}$ decreases) and the fundamental frequency drops: this is why the bass strings of guitars, pianos and the mvet are always the thickest ones.

\textbf{6. Polarization demonstration.}\par
(a) \emph{Rope and slits.} A rope is passed through two vertical slits and shaken vertically: the transverse vibration passes through both slits freely. If the second slit is rotated until it is horizontal (at $90^{\circ}$ to the first), the vertical vibration of the rope cannot pass through it: beyond the second slit the rope does not move. The first slit has \emph{plane polarized} the wave; the second slit acts as an \emph{analyser}.\par
(b) \emph{Polaroid sheets.} The first polaroid sheet plays the role of the first slit: it transmits only the component of the light's vibration parallel to its transmission axis, producing plane polarized light. The second sheet (the analyser) transmits only the component parallel to its own axis. When the two axes are crossed at $90^{\circ}$, no component of the polarized light can pass: the transmitted intensity falls practically to zero and the view becomes dark.\par
(c) \emph{Conclusion.} Only a \emph{transverse} wave can be polarized, because only then does the vibration have a direction in space, perpendicular to the direction of propagation, that can be selected by a slit or a polaroid. The darkening of crossed polaroids therefore proves that \textbf{light is a transverse wave}. Sound in air is a longitudinal wave: its vibration is along the direction of propagation, so no rotation of any ``analyser'' can extinguish it, and polarization of sound in air is impossible.

\textbf{7. Synthesis sentence.}\par
\emph{``Whether it is the fibre of a mvet vibrating at $220\ \mathrm{Hz}$ or a beam of sunlight crossing two polaroid sheets, the physics is the same: the waves are transverse — their speed on a string is fixed by $v=\sqrt{T/\mu}$, and their direction of vibration can be selected by polarization.''}

\begin{aretenir}[What this activity consolidates]
\begin{itemize}
\item Fundamental mode of a string fixed at both ends: $\lambda_1=2L$, so $f_1=\dfrac{1}{2L}\sqrt{\dfrac{T}{\mu}}$.
\item $f_1$ increases as $\sqrt{T}$ and decreases as $1/\sqrt{\mu}$: tighten to sharpen the note, thicken to lower it.
\item Polarization selects the direction of vibration; it is the signature of transverse waves, for mechanical waves on ropes as for light.
\end{itemize}
\end{aretenir}

\newpage

\coursec{Methods and worked examples}

\courssub{Methods and Worked Examples}

\begin{methode}[Identifying a Wave as Transverse]
\begin{itemize}
\item \textbf{Step 1}: Identify the wave motion: transverse waves have oscillations perpendicular to the direction of propagation.
\item \textbf{Step 2}: Determine the medium and the type of wave: transverse waves can be polarized; longitudinal waves cannot.
\item \textbf{Step 3}: Recall the key idea: a transverse wave is polarized when the vibrations are confined to a plane.
\item \textbf{Step 4}: Use the concept of a polarizer: only the component of the wave oscillating in the plane of the polarizer is transmitted.
\end{itemize}
\end{methode}

\begin{exoresolu}[Polarization of a Wave]
\textbf{Statement:} A wave is polarized when the vibrations of the wave are confined to a plane. The medium is a transverse wave. The medium is a transverse wave. The medium is a transverse wave. The medium is a transverse wave. The medium is a transverse wave. The medium is a transverse wave. The medium is a transverse wave. The medium is a transverse wave. The medium is a transverse wave. The medium is a transverse wave. The medium is a transverse wave. The medium is a transverse wave. The medium is a transverse wave. The medium is a transverse wave. The medium is a transverse wave. The medium is a transverse wave. The medium is a transverse wave. The medium is a transverse wave. The medium is a transverse wave. The medium is a transverse wave. The medium is a transverse wave. The medium is a transverse wave. The medium is a transverse wave. The medium is a transverse wave. The medium is a transverse wave. The medium is a transverse wave. The medium is a transverse wave.
\tcblower
Detailed solution to be added here.
\end{exoresolu}

\newpage

\coursec{Exercises}

\courssub{I check my knowledge}

\begin{exercice}[Multiple choice questions]
For each question, choose the single correct answer and justify your choice in one sentence.
\begin{enumerate}
\item Polarization is a phenomenon that can occur:
\begin{enumerate}
\item for all waves, transverse and longitudinal;
\item only for transverse waves;
\item only for longitudinal waves;
\item only for light waves.
\end{enumerate}
\item A plane polarized wave is one in which:
\begin{enumerate}
\item the amplitude decreases with distance;
\item the vibrations of the medium are confined to a single plane containing the direction of propagation;
\item the wave travels in a plane;
\item the frequency is constant.
\end{enumerate}
\item The speed of a transverse wave on a taut string of linear density $\mu$ under tension $T$ is:
\begin{enumerate}
\item $v = \sqrt{T/\mu}$;
\item $v = \sqrt{\mu/T}$;
\item $v = T/\mu$;
\item $v = \sqrt{T\mu}$.
\end{enumerate}
\item Unpolarized light of intensity $I_0$ passes through an ideal polarizer. The transmitted intensity is:
\begin{enumerate}
\item $I_0$;
\item $I_0/2$;
\item $I_0/4$;
\item $0$.
\end{enumerate}
\end{enumerate}
\end{exercice}

\begin{exercice}[True or false]
State whether each statement is true or false, justifying your answer briefly.
\begin{enumerate}
\item Sound waves in air can be polarized by passing them through a narrow slit.
\item Rotating a second polaroid (analyser) through $90^{\circ}$ from the position of maximum transmission reduces the transmitted intensity to zero.
\item Doubling the tension of a stretched wire doubles the speed of transverse waves on it.
\item For a string fixed at both ends, the fundamental wavelength equals twice the length of the string.
\end{enumerate}
\end{exercice}

\begin{exercice}[Definitions]
Define each of the following terms clearly, in one or two sentences each:
\begin{enumerate}
\item plane polarized wave;
\item linear density (mass per unit length) of a string;
\item polarizer and analyser;
\item fundamental mode of vibration of a stretched string.
\end{enumerate}
\end{exercice}

\begin{exercice}[Direct calculations]
\begin{enumerate}
\item A wire of length $2.0\ \mathrm{m}$ has a mass of $5.0\ \mathrm{g}$. Calculate its linear density $\mu$.
\item A transverse wave travels at $80\ \mathrm{m\,s^{-1}}$ along a string under a tension of $64\ \mathrm{N}$. Calculate the linear density of the string.
\item Unpolarized light of intensity $I_0$ falls on a polaroid sheet. A second polaroid has its axis at $30^{\circ}$ to the first. Calculate the final transmitted intensity as a fraction of $I_0$.
\end{enumerate}
\end{exercice}

\courssub{I practise}

\begin{exercice}[Malus's law]
Unpolarized light of intensity $I_0$ passes through two ideal polaroids whose transmission axes make an angle of $60^{\circ}$ with each other.
\begin{enumerate}
\item State Malus's law.
\item Calculate the intensity transmitted by the first polaroid.
\item Calculate the intensity transmitted by the second polaroid as a fraction of $I_0$.
\item What angle between the axes would reduce the final intensity to $I_0/8$?
\end{enumerate}
\end{exercice}

\begin{exercice}[Wave on a steel wire]
A steel wire of length $1.5\ \mathrm{m}$ and mass $12\ \mathrm{g}$ is stretched under a tension of $200\ \mathrm{N}$.
\begin{enumerate}
\item Calculate the linear density $\mu$ of the wire.
\item Calculate the speed $v$ of transverse waves on the wire.
\item A transverse wave of frequency $100\ \mathrm{Hz}$ travels along the wire. Calculate its wavelength.
\item The tension is increased to $800\ \mathrm{N}$. Calculate the new wave speed and comment on the effect of quadrupling the tension.
\end{enumerate}
\end{exercice}

\begin{exercice}[Vibrating string fixed at both ends]
A guitar string of length $0.65\ \mathrm{m}$ and linear density $1.2\times10^{-3}\ \mathrm{kg\,m^{-1}}$ is under a tension of $75\ \mathrm{N}$.
\begin{enumerate}
\item Calculate the speed of transverse waves on the string.
\item Calculate the fundamental frequency of the string.
\item Calculate the frequencies of the second and third harmonics.
\end{enumerate}
\end{exercice}

\begin{exercice}[Changing the string]
A musician replaces a string by another string of double the linear density, keeping the same tension and the same vibrating length.
\begin{enumerate}
\item Show that the fundamental frequency is divided by $\sqrt{2}$.
\item Comment on the effect this has on the pitch of the note heard.
\item What change in tension (at constant length and linear density) would restore the original fundamental frequency?
\end{enumerate}
\end{exercice}

\courssub{I prepare for the GCE}

\begin{exercice}[Polarization of a rope wave and of sound]
\begin{enumerate}
\item Define a plane polarized wave. \hfill (2 marks)
\item Describe, with the aid of a labelled diagram, an experiment using a rope and two slits to demonstrate the polarization of a transverse wave. State clearly what is observed when the second slit is rotated through $90^{\circ}$. \hfill (5 marks)
\item Explain why the result of this experiment proves that the wave on the rope is transverse. \hfill (2 marks)
\item Explain why a sound wave in air cannot be polarized by the same arrangement. \hfill (2 marks)
\item State two applications of polarized light in everyday life or technology. \hfill (2 marks)
\end{enumerate}
\end{exercice}

\begin{exercice}[Sonometer experiment]
In an experiment to study transverse waves on a stretched wire, a sonometer wire is kept under tension by hanging masses. For each value of the tension $T$, the wavelength $\lambda$ of the wave produced at a fixed frequency is measured. The results are:

\begin{center}
\begin{tabular}{|c|c|c|c|c|c|}
\hline
\rowcolor{popblueL} $T$ (N) & 25 & 50 & 75 & 100 & 125 \\
\hline
$\lambda$ (m) & 0.40 & 0.57 & 0.69 & 0.80 & 0.89 \\
\hline
\end{tabular}
\end{center}

The frequency of the source is $f = 256\ \mathrm{Hz}$.
\begin{enumerate}
\item Write down the relation connecting the wave speed $v$, the tension $T$ and the linear density $\mu$ of the wire. \hfill (1 mark)
\item Show that $\lambda^2 = \dfrac{T}{\mu f^2}$. \hfill (2 marks)
\item Copy and complete the table by computing $\lambda^2$ for each measurement. \hfill (2 marks)
\item Plot a graph of $\lambda^2$ (vertical axis) against $T$ (horizontal axis). \hfill (4 marks)
\item Determine the gradient of the graph. \hfill (2 marks)
\item Deduce the linear density $\mu$ of the wire. \hfill (2 marks)
\item State one precaution you would take to obtain accurate results. \hfill (1 mark)
\end{enumerate}
\end{exercice}

\begin{exercice}[Guitar string in its fundamental mode]
A guitar string fixed at both ends has length $L = 0.80\ \mathrm{m}$ and linear density $\mu = 4.0\times10^{-3}\ \mathrm{kg\,m^{-1}}$. It is stretched under a tension $T = 128\ \mathrm{N}$ and vibrates in its fundamental mode.
\begin{enumerate}
\item Draw a labelled diagram of the string vibrating in its fundamental mode, indicating nodes and antinodes. \hfill (2 marks)
\item Show that the fundamental frequency is given by $f_1 = \dfrac{1}{2L}\sqrt{\dfrac{T}{\mu}}$. \hfill (3 marks)
\item Calculate the fundamental frequency of the string. \hfill (2 marks)
\item Calculate the new tension required to raise the fundamental frequency by a factor of $1.5$, the length and linear density being unchanged. \hfill (3 marks)
\item State two assumptions made in deriving the formula used in (b). \hfill (2 marks)
\item Explain, in terms of wave speed and wavelength, why increasing the tension raises the pitch of the note. \hfill (2 marks)
\end{enumerate}
\end{exercice}



\newpage

\coursec{Solutions to the exercises}

\courssub{I check my knowledge}

\begin{corrige}[Exercise 1 --- Multiple choice questions]
\begin{enumerate}
\item \textbf{(b)}. Polarization requires the vibration to have a direction perpendicular to the propagation; longitudinal waves vibrate along the propagation direction and cannot be polarized.
\item \textbf{(b)}. This is precisely the definition of plane (linear) polarization: the vibrations of the medium are confined to a single plane containing the direction of propagation.
\item \textbf{(a)}. $v=\sqrt{T/\mu}$; dimensional check: $[T/\mu]=\mathrm{N\,kg^{-1}\,m}=\mathrm{kg\,m\,s^{-2}\,kg^{-1}\,m}=\mathrm{m^2\,s^{-2}}$, whose square root is a speed.
\item \textbf{(b)}. An ideal polarizer transmits only the component of the vibration parallel to its axis; averaging over all random orientations of the unpolarized light gives $I_0/2$.
\end{enumerate}
\end{corrige}

\begin{corrige}[Exercise 2 --- True or false]
\begin{enumerate}
\item \textbf{False.} Sound in air is longitudinal; its vibrations are along the direction of propagation, so no slit orientation can select a vibration direction.
\item \textbf{True.} By Malus's law $I=I_{\max}\cos^2\theta$; at $\theta=90^{\circ}$, $\cos\theta=0$ and $I=0$ (crossed polaroids).
\item \textbf{False.} Since $v=\sqrt{T/\mu}$, doubling $T$ multiplies $v$ by $\sqrt{2}\approx 1.41$, not by $2$.
\item \textbf{True.} In the fundamental mode the string carries one loop, so $L=\lambda_1/2$, i.e. $\lambda_1=2L$.
\end{enumerate}
\end{corrige}

\begin{corrige}[Exercise 3 --- Definitions]
\begin{enumerate}
\item A \cle{plane polarized wave} is a transverse wave in which the vibrations of the medium are confined to one fixed plane containing the direction of propagation.
\item The \cle{linear density} $\mu$ of a string is its mass per unit length: $\mu=m/L$, in $\mathrm{kg\,m^{-1}}$.
\item A \cle{polarizer} is a device that produces plane polarized light from unpolarized light; an \cle{analyser} is a second identical device used to detect the polarization and measure the orientation of its plane (by Malus's law).
\item The \cle{fundamental mode} is the simplest standing-wave pattern of a string fixed at both ends: a single loop with nodes at the ends and one antinode at the centre, of wavelength $\lambda_1=2L$ and frequency $f_1=v/(2L)$.
\end{enumerate}
\end{corrige}

\begin{corrige}[Exercise 4 --- Direct calculations]
\begin{enumerate}
\item $\mu=\dfrac{m}{L}=\dfrac{5.0\times10^{-3}}{2.0}=2.5\times10^{-3}\ \mathrm{kg\,m^{-1}}$.
\item From $v=\sqrt{T/\mu}$ we get $\mu=\dfrac{T}{v^2}=\dfrac{64}{80^2}=\dfrac{64}{6400}=1.0\times10^{-2}\ \mathrm{kg\,m^{-1}}$.
\item After the first polaroid: $I_1=I_0/2$. After the second (Malus's law):
\[I_2=I_1\cos^2 30^{\circ}=\frac{I_0}{2}\times\left(\frac{\sqrt{3}}{2}\right)^2=\frac{I_0}{2}\times\frac{3}{4}=\frac{3I_0}{8}=0.375\,I_0.\]
\end{enumerate}
\end{corrige}

\courssub{I practise}

\begin{corrige}[Exercise 5 --- Malus's law]
\begin{enumerate}
\item \textbf{Malus's law:} when plane polarized light of intensity $I_1$ falls on an ideal analyser whose axis makes an angle $\theta$ with the plane of polarization, the transmitted intensity is $I=I_1\cos^2\theta$.
\item The first polaroid receives unpolarized light, so $I_1=I_0/2$.
\item $I_2=I_1\cos^2 60^{\circ}=\dfrac{I_0}{2}\times(0.5)^2=\dfrac{I_0}{8}=0.125\,I_0$.
\item We require $\dfrac{I_0}{2}\cos^2\theta=\dfrac{I_0}{8}$, so $\cos^2\theta=\dfrac{1}{4}$, $\cos\theta=\dfrac{1}{2}$, hence $\theta=60^{\circ}$.
\end{enumerate}
\end{corrige}

\begin{corrige}[Exercise 6 --- Wave on a steel wire]
\begin{enumerate}
\item $\mu=\dfrac{m}{L}=\dfrac{12\times10^{-3}}{1.5}=8.0\times10^{-3}\ \mathrm{kg\,m^{-1}}$.
\item $v=\sqrt{\dfrac{T}{\mu}}=\sqrt{\dfrac{200}{8.0\times10^{-3}}}=\sqrt{2.5\times10^{4}}\approx 1.6\times10^{2}\ \mathrm{m\,s^{-1}}$ ($158\ \mathrm{m\,s^{-1}}$).
\item $\lambda=\dfrac{v}{f}=\dfrac{158}{100}\approx 1.6\ \mathrm{m}$.
\item $v'=\sqrt{\dfrac{800}{8.0\times10^{-3}}}=\sqrt{1.0\times10^{5}}\approx 3.2\times10^{2}\ \mathrm{m\,s^{-1}}$ ($316\ \mathrm{m\,s^{-1}}$). Quadrupling the tension only \emph{doubles} the speed, because $v\propto\sqrt{T}$.
\end{enumerate}
\end{corrige}

\begin{corrige}[Exercise 7 --- Vibrating string fixed at both ends]
\begin{enumerate}
\item $v=\sqrt{\dfrac{T}{\mu}}=\sqrt{\dfrac{75}{1.2\times10^{-3}}}=\sqrt{6.25\times10^{4}}=2.5\times10^{2}\ \mathrm{m\,s^{-1}}$.
\item $f_1=\dfrac{v}{2L}=\dfrac{250}{2\times0.65}=\dfrac{250}{1.30}\approx 1.9\times10^{2}\ \mathrm{Hz}$ ($192\ \mathrm{Hz}$).
\item $f_2=2f_1\approx 3.8\times10^{2}\ \mathrm{Hz}$ ($385\ \mathrm{Hz}$); $f_3=3f_1\approx 5.8\times10^{2}\ \mathrm{Hz}$ ($577\ \mathrm{Hz}$).
\end{enumerate}
\end{corrige}

\begin{corrige}[Exercise 8 --- Changing the string]
\begin{enumerate}
\item $f_1=\dfrac{1}{2L}\sqrt{\dfrac{T}{\mu}}$. With $\mu'=2\mu$:
\[f_1'=\frac{1}{2L}\sqrt{\frac{T}{2\mu}}=\frac{1}{\sqrt{2}}\times\frac{1}{2L}\sqrt{\frac{T}{\mu}}=\frac{f_1}{\sqrt{2}}.\]
\item The frequency decreases, so the note heard is \emph{lower} in pitch (a heavier string sounds deeper, as with the bass strings of a guitar or piano).
\item Since $f_1\propto\sqrt{T}$ at constant $L$ and $\mu$, restoring $f_1$ while $\mu$ is doubled requires $\sqrt{T'/(2\mu)}=\sqrt{T/\mu}$, i.e. $T'=2T$: the tension must be \emph{doubled}.
\end{enumerate}
\end{corrige}

\courssub{I prepare for the GCE}

\begin{corrige}[Exercise 9 --- Polarization of a rope wave and of sound]
\begin{enumerate}
\item A plane polarized wave is a transverse wave in which the vibrations of the medium are confined to a single fixed plane containing the direction of propagation. \hfill (2 marks: \emph{transverse} + \emph{single plane containing the propagation direction})
\item A long rope is fixed at one end and shaken at the other end so that it vibrates in all transverse directions (unpolarized wave). The rope passes through a vertical slit $S_1$ in a board: only the vertical component of the vibration is transmitted beyond $S_1$, so the wave between $S_1$ and a second slit $S_2$ is plane polarized in the vertical plane. When $S_2$ is also vertical, the wave passes through with almost unchanged amplitude. When $S_2$ is rotated through $90^{\circ}$ (horizontal slit), the vertical vibration cannot pass: the rope beyond $S_2$ remains at rest --- the wave is completely stopped.
\begin{center}
\begin{tikzpicture}[scale=0.9]
\draw[thick,popblue] (-3,0) sin (-2.5,0.4) cos (-2,0) sin (-1.5,-0.4) cos (-1,0);
\draw[thick,popdark] (-0.6,-1) -- (-0.6,-0.35) (-0.6,0.35) -- (-0.6,1);
\draw[thick,popdark] (-0.3,-1) -- (-0.3,-0.35) (-0.3,0.35) -- (-0.3,1);
\node[below] at (-0.45,-1) {$S_1$};
\draw[thick,popblue] (0,0) sin (0.5,0.4) cos (1,0) sin (1.5,-0.4) cos (2,0);
\draw[thick,popdark] (2.4,-1) -- (2.4,-0.35) (2.4,0.35) -- (2.4,1);
\draw[thick,popdark] (2.7,-1) -- (2.7,-0.35) (2.7,0.35) -- (2.7,1);
\node[below] at (2.55,-1) {$S_2$};
\draw[thick,popblue] (3,0) sin (3.5,0.4) cos (4,0) sin (4.5,-0.4) cos (5,0);
\node[left] at (-3,0) {hand};
\node[right] at (5,0) {fixed end};
\draw[->,thick,poporange] (-2.9,0.8) -- (-1.9,0.8) node[midway,above] {propagation};
\end{tikzpicture}
\end{center}
\hfill (5 marks: labelled diagram 2; role of $S_1$ as polarizer 1; observation with $S_2$ parallel 1; observation with $S_2$ at $90^{\circ}$ 1)
\item The wave is stopped when the slit is perpendicular to the vibration plane, which shows the vibration has a definite direction \emph{perpendicular} to the propagation; only transverse waves possess such a transverse vibration direction that a slit can select or block. \hfill (2 marks)
\item Sound in air is longitudinal: the air layers vibrate \emph{along} the direction of propagation. A slit of any orientation is equally traversed, so no orientation can block the wave; hence sound cannot be polarized. \hfill (2 marks)
\item Applications (any two): polarizing (Polaroid) sunglasses to reduce glare reflected by roads and water; polarizing filters in photography; liquid-crystal displays (LCD screens); stress analysis in transparent materials (photoelasticity). \hfill (2 marks)
\end{enumerate}
\textbf{Examiner's expectations:} the diagram must show the rope, both slits and the direction of propagation; the key marking point in (c) is the explicit link ``blocked wave $\Rightarrow$ vibration perpendicular to propagation''.
\end{corrige}

\begin{corrige}[Exercise 10 --- Sonometer experiment]
\begin{enumerate}
\item $v=\sqrt{\dfrac{T}{\mu}}$. \hfill (1 mark)
\item Since $v=f\lambda$, we have $f\lambda=\sqrt{T/\mu}$. Squaring: $f^2\lambda^2=\dfrac{T}{\mu}$, hence $\lambda^2=\dfrac{T}{\mu f^2}$. \hfill (2 marks)
\item Completed table:
\begin{center}
\begin{tabular}{|c|c|c|c|c|c|}
\hline
\rowcolor{popblueL} $T$ (N) & 25 & 50 & 75 & 100 & 125 \\
\hline
$\lambda$ (m) & 0.40 & 0.57 & 0.69 & 0.80 & 0.89 \\
\hline
$\lambda^2$ ($\mathrm{m^2}$) & 0.160 & 0.325 & 0.476 & 0.640 & 0.792 \\
\hline
\end{tabular}
\end{center}
\hfill (2 marks: all five values correct)
\item Graph: axes labelled with quantities and units, suitable uniform scales, points plotted correctly, best straight line drawn through the origin.
\begin{center}
\begin{tikzpicture}
\begin{axis}[width=0.75\linewidth,height=6cm,
xlabel={$T$ (N)},ylabel={$\lambda^2$ ($\mathrm{m^2}$)},
xmin=0,xmax=140,ymin=0,ymax=0.9,
grid=both,grid style={popline,very thin},
tick label style={font=\small},label style={font=\small}]
\addplot[only marks,mark=*,popblue] coordinates {(25,0.160) (50,0.325) (75,0.476) (100,0.640) (125,0.792)};
\addplot[domain=0:140,poporange,thick] {0.00632*x};
\end{axis}
\end{tikzpicture}
\end{center}
The points lie on a straight line through the origin, confirming $\lambda^2\propto T$. \hfill (4 marks: axes + scales 1; points 2; line through origin 1)
\item Gradient, using the points $(25;\,0.160)$ and $(125;\,0.792)$ on the line:
\[k=\frac{0.792-0.160}{125-25}=\frac{0.632}{100}=6.3\times10^{-3}\ \mathrm{m^2\,N^{-1}}.\]
\hfill (2 marks: method 1; value with unit 1)
\item Since $k=\dfrac{1}{\mu f^2}$, we get
\[\mu=\frac{1}{k f^2}=\frac{1}{6.3\times10^{-3}\times 256^2}=\frac{1}{6.3\times10^{-3}\times 6.55\times10^{4}}\approx 2.4\times10^{-3}\ \mathrm{kg\,m^{-1}}.\]
\hfill (2 marks: relation 1; value with unit 1)
\item Precaution (any one): ensure the pulley is nearly frictionless so the tension equals the weight of the hanging masses; measure the wavelength over several loops and divide; avoid parallax when reading the positions of the bridges; keep the frequency of the vibrator constant. \hfill (1 mark)
\end{enumerate}
\textbf{Examiner's expectations:} the gradient must be taken from the drawn line (not from a single data pair) over at least half the line length; the final value of $\mu$ must carry the unit $\mathrm{kg\,m^{-1}}$.
\end{corrige}

\begin{corrige}[Exercise 11 --- Guitar string in its fundamental mode]
\begin{enumerate}
\item Diagram: a single loop between the two fixed ends; nodes $N$ at both ends, one antinode $A$ at the centre.
\begin{center}
\begin{tikzpicture}[scale=1.0]
\draw[thick] (0,0) -- (6,0);
\draw[thick,popblue] (0,0) .. controls (1.5,1.2) and (4.5,1.2) .. (6,0);
\draw[dashed,popmuted] (0,0) .. controls (1.5,-1.2) and (4.5,-1.2) .. (6,0);
\fill[popdark] (0,0) circle (0.07);
\fill[popdark] (6,0) circle (0.07);
\fill[poporange] (3,0.95) circle (0.07);
\node[below] at (0,-0.1) {$N$};
\node[below] at (6,-0.1) {$N$};
\node[above] at (3,1.0) {$A$};
\draw[<->] (0,-0.7) -- (6,-0.7) node[midway,below] {$L=\lambda_1/2$};
\end{tikzpicture}
\end{center}
\hfill (2 marks: single loop 1; nodes and antinode labelled 1)
\item In the fundamental mode the string carries one loop, so $L=\dfrac{\lambda_1}{2}$, i.e. $\lambda_1=2L$. The wave speed on the string is $v=\sqrt{T/\mu}$, and $v=f_1\lambda_1$. Hence
\[f_1=\frac{v}{\lambda_1}=\frac{1}{2L}\sqrt{\frac{T}{\mu}}.\]
\hfill (3 marks: $\lambda_1=2L$ 1; $v=\sqrt{T/\mu}$ 1; combination 1)
\item $f_1=\dfrac{1}{2\times0.80}\sqrt{\dfrac{128}{4.0\times10^{-3}}}=\dfrac{1}{1.60}\sqrt{3.2\times10^{4}}=0.625\times 178.9\approx 1.1\times10^{2}\ \mathrm{Hz}$ ($112\ \mathrm{Hz}$). \hfill (2 marks)
\item Since $f_1\propto\sqrt{T}$ at constant $L$ and $\mu$, raising $f_1$ by a factor $1.5$ requires multiplying $T$ by $1.5^2=2.25$:
\[T'=2.25\times 128=288\ \mathrm{N}.\]
\hfill (3 marks: proportionality 1; squaring the factor 1; value 1)
\item Assumptions (any two): the string is perfectly flexible (no stiffness); its linear density is uniform; the tension is the same at every point and unchanged by the small vibrations; the amplitude of vibration is small; energy losses (damping) are neglected. \hfill (2 marks)
\item Increasing $T$ increases the wave speed $v=\sqrt{T/\mu}$. The vibrating length fixes the fundamental wavelength ($\lambda_1=2L$), so $f_1=v/\lambda_1$ increases with $v$. A higher frequency is perceived as a higher-pitched note. \hfill (2 marks)
\end{enumerate}
\textbf{Examiner's expectations:} in (d) the most common error is to multiply $T$ by $1.5$ instead of $1.5^2$; the square-root dependence must be stated explicitly to earn the method mark.
\end{corrige}

\newpage

\coursec{Summary sheet}

\begin{popfigure}
\begin{center}\tcbox[colback=popink,colframe=popink,coltext=white,arc=9pt,boxsep=4pt,left=6pt,right=6pt]{\bfseries\small Mechanical Waves (continued)}\end{center}
\vspace{-2pt}
\begin{tcbraster}[raster columns=3,raster equal height=rows,raster column skip=6pt,raster row skip=6pt,enhanced,blankest]
\begin{tcolorbox}[enhanced,colback=white,colframe=popblue,boxrule=0.9pt,arc=3pt,title={Polarization: concept},fonttitle=\bfseries\scriptsize,coltitle=white,colbacktitle=popblue,left=4pt,right=4pt,top=2pt,bottom=2pt,boxsep=1pt,toptitle=1.5pt,bottomtitle=1.5pt,halign=left]
\begin{itemize}[nosep,leftmargin=*,itemsep=1pt,font=\scriptsize]
\item Transverse vs longitudinal
\item Slit model, rope demo
\end{itemize}
\end{tcolorbox}
\begin{tcolorbox}[enhanced,colback=white,colframe=poporange,boxrule=0.9pt,arc=3pt,title={Plane polarized light},fonttitle=\bfseries\scriptsize,coltitle=white,colbacktitle=poporange,left=4pt,right=4pt,top=2pt,bottom=2pt,boxsep=1pt,toptitle=1.5pt,bottomtitle=1.5pt,halign=left]
\begin{itemize}[nosep,leftmargin=*,itemsep=1pt,font=\scriptsize]
\item Polaroid, reflection, scattering
\item Analyser: Malus $I=I_0\cos^2\theta$
\end{itemize}
\end{tcolorbox}
\begin{tcolorbox}[enhanced,colback=white,colframe=popgreen,boxrule=0.9pt,arc=3pt,title={Speed on taut string},fonttitle=\bfseries\scriptsize,coltitle=white,colbacktitle=popgreen,left=4pt,right=4pt,top=2pt,bottom=2pt,boxsep=1pt,toptitle=1.5pt,bottomtitle=1.5pt,halign=left]
\begin{itemize}[nosep,leftmargin=*,itemsep=1pt,font=\scriptsize]
\item $v=\sqrt{T/\mu}$
\item $\mu$ = linear density
\end{itemize}
\end{tcolorbox}
\begin{tcolorbox}[enhanced,colback=white,colframe=poppurple,boxrule=0.9pt,arc=3pt,title={Factors affecting $v$},fonttitle=\bfseries\scriptsize,coltitle=white,colbacktitle=poppurple,left=4pt,right=4pt,top=2pt,bottom=2pt,boxsep=1pt,toptitle=1.5pt,bottomtitle=1.5pt,halign=left]
\begin{itemize}[nosep,leftmargin=*,itemsep=1pt,font=\scriptsize]
\item Tension $T$ (hanging mass)
\item Material, thickness
\end{itemize}
\end{tcolorbox}
\begin{tcolorbox}[enhanced,colback=white,colframe=poppink,boxrule=0.9pt,arc=3pt,title={Sonometer},fonttitle=\bfseries\scriptsize,coltitle=white,colbacktitle=poppink,left=4pt,right=4pt,top=2pt,bottom=2pt,boxsep=1pt,toptitle=1.5pt,bottomtitle=1.5pt,halign=left]
\begin{itemize}[nosep,leftmargin=*,itemsep=1pt,font=\scriptsize]
\item Resonance with tuning fork
\item $f=\dfrac{1}{2L}\sqrt{T/\mu}$
\end{itemize}
\end{tcolorbox}
\begin{tcolorbox}[enhanced,colback=white,colframe=popgold,boxrule=0.9pt,arc=3pt,title={Applications},fonttitle=\bfseries\scriptsize,coltitle=white,colbacktitle=popgold,left=4pt,right=4pt,top=2pt,bottom=2pt,boxsep=1pt,toptitle=1.5pt,bottomtitle=1.5pt,halign=left]
\begin{itemize}[nosep,leftmargin=*,itemsep=1pt,font=\scriptsize]
\item Guitar, piano tuning
\item Sunglasses, LCD, stress analysis
\end{itemize}
\end{tcolorbox}
\end{tcbraster}

\legende{Mind map of the chapter: polarization of transverse waves and the speed of transverse waves on taut strings.}
\end{popfigure}

\begin{aretenir}[Summary Sheet --- Mechanical Waves (Continued)]
\textbf{Key definitions}
\begin{itemize}
\item \cle{Polarization}: property of a \textbf{transverse} wave in which the vibrations, which would otherwise occur in all possible directions perpendicular to the direction of propagation, are restricted to a definite pattern of directions. A \cle{plane polarized} wave vibrates in one fixed plane only (the plane of polarization), which contains the direction of propagation.
\item Only \textbf{transverse} waves can be polarized; longitudinal waves (sound in air) cannot, because their vibrations are already along the single direction of propagation. The fact that light can be polarized is the experimental proof that light is a \textbf{transverse} wave.
\item \cle{Polarizer}: device producing plane polarized light from ordinary (unpolarized) light; \cle{analyser}: an identical device used to detect the polarization. When polarizer and analyser are \emph{crossed} (transmission axes at $90^{\circ}$), no light passes: darkness.
\item \cle{Linear density} (mass per unit length) of a string: $\mu = m/L$, SI unit $\mathrm{kg\,m^{-1}}$.
\end{itemize}

\textbf{Laws and formulas to know}
\begin{itemize}
\item Speed of a transverse wave on a taut string:
\[ v=\sqrt{\dfrac{T}{\mu}} \qquad \text{($T$ tension in N, $\mu$ in $\mathrm{kg\,m^{-1}}$, $v$ in $\mathrm{m\,s^{-1}}$)} \]
$v$ increases when the tension increases and decreases when the linear density increases; $v$ does \emph{not} depend on the frequency or the amplitude of the wave.
\item Since $\mu=\rho S$ ($\rho$ density of the material, $S$ cross-sectional area of the wire): $v=\sqrt{T/(\rho S)}$. Hence, at equal tension, a thicker wire (larger $S$) carries slower waves.
\item Fundamental frequency of a string fixed at both ends (sonometer wire), from $v=f\lambda$ with $\lambda_1=2L$:
\[ f_1=\dfrac{1}{2L}\sqrt{\dfrac{T}{\mu}}, \qquad f_n=n f_1 \quad (n=1,2,3,\dots) \]
\item Law of Malus (analyser whose axis makes an angle $\theta$ with that of the polarizer): $I=I_0\cos^2\theta$. Unpolarized light passing through a first polarizer loses half its intensity: $I=I_0/2$.
\end{itemize}

\textbf{Methods}
\begin{itemize}
\item \textbf{Computing $v$ on a string:} convert the mass per metre to $\mu$ in $\mathrm{kg\,m^{-1}}$; if a mass $M$ hangs over a pulley, the tension is $T=Mg$; then apply $v=\sqrt{T/\mu}$.
\item \textbf{Sonometer problems:} identify $L$, $T$ and $\mu$; at resonance with a tuning fork of frequency $f$, write $f=\dfrac{1}{2L}\sqrt{T/\mu}$ and solve for the unknown quantity.
\item \textbf{Polarization questions:} state that since light can be polarized, light must be a transverse wave; support the explanation with the mechanical analogy of a rope passing through two slits (the wave passes only when the slits are parallel).
\end{itemize}

\textbf{Common mistakes}
\begin{itemize}
\item Claiming that sound can be polarized --- false: sound in air is a longitudinal wave.
\item Using $\mu$ in $\mathrm{g\,cm^{-1}}$ or $\mathrm{g\,m^{-1}}$ without converting to $\mathrm{kg\,m^{-1}}$.
\item Forgetting that $T=Mg$ uses $g=9.81\ \mathrm{m\,s^{-2}}$ (or $10\ \mathrm{m\,s^{-2}}$ if the question says so); the tension is a force, not the mass $M$ alone.
\item Confusing the \emph{plane of polarization} (the plane in which the vibration occurs) with the \emph{direction of propagation}: the plane of polarization \emph{contains} the direction of propagation.
\item Believing that $v$ depends on frequency: on a given string under given tension, $v$ is fixed by $T$ and $\mu$; changing $f$ only changes the wavelength $\lambda=v/f$.
\end{itemize}

\textbf{GCE tips}
\begin{itemize}
\item Always justify ``light is a transverse wave'' by describing the experiment with two Polaroids: rotating the analyser until the light is extinguished proves the vibrations are confined to one plane.
\item Show units at each step of a calculation and give $v$ with sensible significant figures (two or three).
\item In sonometer questions, state the condition for resonance (frequency of the vibrating string = frequency of the tuning fork) before substituting numerical values.
\item Quote applications: Polaroid sunglasses (which cut glare from light partially polarized by reflection), LCD screens, photoelastic stress analysis, and the tuning of guitars and pianos (tightening a peg raises $T$, hence raises $f$).
\end{itemize}
\end{aretenir}

\findecours

\end{document}
